\documentclass[oneside,polski,logo]{amuthesis}
\usepackage[font=small,labelfont=bf]{caption}

% Zdefiniuj kodowanie pliku źródłowego (domyślnie utf8)
\usepackage{booktabs}
\usepackage{tikz}
\usepackage{xcolor}
\usetikzlibrary{positioning, arrows.meta, calc}
\usepackage[utf8]{inputenc}
\usepackage{authblk}
\usepackage{multirow}
\usepackage[utf8]{inputenc}
\usepackage[T1]{fontenc}
\usepackage[polish]{babel}
\usepackage{array}
\usepackage{caption}
\usepackage{graphicx}
\usepackage{lipsum}
\usepackage[utf8]{inputenc}
\usepackage[polish]{babel}
\usepackage{csquotes}
\usepackage[backend=biber, style=ieee, sorting=nyt, bibencoding=utf8]{biblatex}
\addbibresource{bibliografia.bib}
\usepackage[compatibility=false]{caption}
\usepackage{subcaption}


% ======================================================== %
% Dane autorów i pracy                                      %
% ======================================================== %

% --- Autorzy prac
% --- --- Imię i nazwisko
% --- --- Numer albumu
% --- --- Płeć (M/K)
\firstAuthor{Wiktoria Grzesik}
\firstAlbum{473561}
\stsexFirstAuthor{K}

%\secondAuthor{Autor Drugi}
%\secondAlbum{234561}
%\stsexSecondAuthor{M}

%\thirdAuthor{Autor Trzeci}
%\thirdAlbum{345612}
%\stsexThirdAuthor{M}

%\fourthAuthor{Autor Czwarty}
%\fourthAlbum{456123}
%\stsexFourthAuthor{F}

% --- Tytuł pracy (w języku polskim i angielskim)
\titlePL{Analiza porównawcza metod uczenia maszynowego i głębokiego w wykrywaniu fake newsów}
\titleEN{Comparative Analysis of Machine Learning and Deep Learning Methods for Fake News Detection}

% --- Typ pracy (inżynierska, licencjacka, magisterska)
\type{magisterska}
% --- Wydział (wykaz skrótów):
% --- --- WA    --- Wydział Anglistyki
% --- --- WAiK  --- Wydział Antropologii i Kulturoznawstwa
% --- --- WAr   --- Wydział Archeologii
% --- --- WB    --- Wydział Biologii
% --- --- WCh   --- Wydział Chemii
% --- --- WFPiK --- Wydział Filologii Polskiej i Klasycznej
% --- --- WFi   --- Wydział Filozoficzny
% --- --- WF    --- Wydział Fizyki
% --- --- WGSE  --- Wydział Geografii Społeczno-Ekonomicznej i Gospodarki Przestrzennej
% --- --- WH    --- Wydział Historii
% --- --- WMiI  --- Wydział Matematyki i Informatyki
% --- --- WNGiG --- Wydział Nauk Geograficznych i Geologicznych
% --- --- WNoS  --- Wydział Nauk o Sztuce
% --- --- WNPiD --- Wydział Nauk Politycznych i Dziennikarstwa
% --- --- WN    --- Wydział Neofilologii
% --- --- WPiK  --- Wydział Psychologii i Kognitywistyki
% --- --- WPAK  --- Wydział Pedagogiczno-Artystyczny w Kaliszu
% --- --- WPiA  --- Wydział Prawa i Administracji
% --- --- WS    --- Wydział Socjologii
% --- --- WSE   --- Wydział Studiów Edukacyjnych
% --- --- WT    --- Wydział Teologiczny
% --- --- IKE   --- Instytut Kultury Europejskiej w Gnieźnie
\faculty{WMiI}
% --- Kierunek (w mianowniku)
\field{informatyka}
% --- Specjalność (w formie mianownikowej)
% --- (ustaw puste, jeśli bez specjalności)
\specialty{Sztuczna Inteligencja}
% --- Promotor (w dopełniaczu)
\supervisor{dra Rafała Jaworskiego}
% --- Data złożenia pracy (Miasto, miesiąc rok)
\date{Poznań, czerwiec 2026}

% --- Zgoda na udostępnienie pracy w czytelni (TAK/NIE)
\stread{TAK}
% --- Zgoda na udostępnienie pracy w zakresie ochrony (TAK/NIE)
\stprotect{TAK}
% --- Data podpisania oświadczenia (Miasto, data)
\stdate{Poznań, \today{} r.}

% ======================================================== %
% Dodatkowe pakiety wykorzystywane w pracy                 %
% ======================================================== %
% ======================================================== %
% Zasadnicza część dokumentu                               %
% ======================================================== %

\begin{document}

% Strona tytułowa
\maketitle
% Oświadczenie
\makestatement

% Blok abstraktu w języku polskim
\begin{streszczenie}
\lipsum[1]

\paragraph{Słowa kluczowe:} klasa
\end{streszczenie}

% Blok abstraktu w języku angielskim
\begin{abstract}
\lipsum[2]

\paragraph{Keywords:} klasa
\end{abstract}

% Opcjonalny blok dedykacji
\begin{dedykacja}
Tu możesz umieścić swoją dedykację.
\end{dedykacja}

% Spis treści
\tableofcontents

% ======================================================== %
% Właściwa część pracy                                     %
% ======================================================== %

\chapter{Wstęp}

\section{Motywacja}
Powszechny dostęp do Internetu i mediów społecznościowych radykalnie przyspieszył obieg informacji, stwarzając jednocześnie warunki sprzyjające masowemu rozpowszechnianiu treści nieprawdziwych i manipulacyjnych, określanych mianem \textit{fake newsów}. Dezinformacja stanowi poważny problem społeczny, polityczny i technologiczny: jak wykazano w~\cite{Shu2017}, fałszywe informacje mogą kształtować opinię publiczną i wpływać na decyzje polityczne, a~badania Vosoughiego i~in.\ \cite{vosoughi2018spread} dowiodły, że dezinformacja rozprzestrzenia się w~mediach społecznościowych szybciej i~osiąga większy zasięg niż przekazy rzetelne. Stanowi to zagrożenie dla jakości debaty publicznej oraz zaufania społecznego.

Problem ulega dalszemu pogłębieniu wskutek rozwoju dużych modeli językowych (\textit{Large Language Models}, LLM), które umożliwiają automatyczne generowanie realistycznych treści naśladujących styl profesjonalnych artykułów informacyjnych~\cite{Li2024}. Jak wskazują Yi i~Huang~\cite{Yi2025}, architektura transformerów -- leżąca u~podstaw nowoczesnych LLM -- stwarza nowe wyzwania dla systemów detekcji, ponieważ syntetycznie wygenerowany tekst staje się coraz trudniejszy do odróżnienia od autentycznych wiadomości.

W obliczu lawinowo rosnącej liczby publikowanych treści ręczna weryfikacja informacji jest niewystarczalna. Rośnie zatem znaczenie automatycznych metod wykrywania dezinformacji opartych na uczeniu maszynowym, uczeniu głębokim oraz przetwarzaniu języka naturalnego (NLP). Współczesne systemy detekcji fake newsów sięgają po całe spektrum podejść -- od klasycznych algorytmów statystycznych po zaawansowane modele transformerowe~\cite{Alshuwaier2025}.

Motywacją do podjęcia niniejszej pracy jest kompleksowa ocena skuteczności wybranych metod wykrywania fake newsów oraz analiza kompromisu pomiędzy jakością klasyfikacji a~wydajnością obliczeniową. Szczególną uwagę poświęcono porównaniu klasycznych modeli uczenia maszynowego, sieci głębokich i~modeli transformerowych w~kontekście ich zastosowania w~systemach działających w~czasie rzeczywistym.

\section{Cel i zakres pracy}
Głównym celem niniejszej pracy jest analiza porównawcza metod uczenia maszynowego, uczenia głębokiego i~modeli transformerowych w~zadaniu wykrywania fake newsów. W~ramach tego celu sformułowano następujące cele szczegółowe:
\begin{itemize} 
    \item analiza zjawiska dezinformacji i~przegląd współczesnych metod jej detekcji,
    \item przygotowanie i~charakterystyka zbiorów danych przeznaczonych do klasyfikacji fake newsów,
    \item implementacja klasycznych algorytmów uczenia maszynowego,
    \item implementacja sieci głębokich oraz modeli transformerowych,
    \item porównanie skuteczności modeli przy użyciu standardowych metryk ewaluacyjnych,
    \item analiza kompromisu pomiędzy jakością klasyfikacji a~czasem inferencji,
    \item implementacja rozszerzenia do przeglądarki internetowej służącego do oceny wiarygodności tekstu.
\end{itemize}

\section{Pytania badawcze}
W pracy postawiono następujące pytania badawcze:
\begin{enumerate} 
    \item Czy modele transformerowe osiągają wyższą skuteczność klasyfikacji fake newsów niż klasyczne algorytmy uczenia maszynowego i~sieci głębokie?
    \item W jaki sposób metoda reprezentacji tekstu wpływa na skuteczność detekcji fake newsów?
    \item Jaki jest kompromis pomiędzy jakością klasyfikacji a~czasem inferencji poszczególnych modeli?
    \item Czy lżejsze modele uczenia maszynowego mogą stanowić efektywną alternatywę dla transformerów w~zastosowaniach wymagających niskich opóźnień?
\end{enumerate}

\section{Hipoteza badawcza}
W pracy przyjęto hipotezę, że modele transformerowe osiągają najwyższą skuteczność w~detekcji fake newsów dzięki zdolności do modelowania długodystansowych zależności kontekstowych i~semantycznych. Jednocześnie zakłada się, że ich przewaga klasyfikacyjna wiąże się z~wyraźnie wyższą złożonością obliczeniową i~dłuższym czasem inferencji w~porównaniu z~klasycznymi metodami uczenia maszynowego. Pomocniczo zakłada się, że odpowiednio dobrane klasyczne algorytmy mogą stanowić efektywną alternatywę w~scenariuszach wymagających niskich opóźnień i~ograniczonych zasobów obliczeniowych.

\section{Struktura pracy}
Praca składa się z~sześciu rozdziałów. Niniejszy rozdział wprowadzający określa cel, zakres i~hipotezę badawczą. Rozdział drugi stanowi przegląd literatury: omówiono w~nim zjawisko dezinformacji, ewolucję metod automatycznej detekcji fake newsów oraz zidentyfikowane luki badawcze. Rozdział trzeci opisuje metodologię: charakterystykę zbiorów danych, proces przygotowania danych, metody reprezentacji tekstu, analizowane modele oraz zestaw metryk ewaluacyjnych. Rozdział czwarty prezentuje wyniki eksperymentów i~analizę porównawczą skuteczności modeli z~uwzględnieniem czasu inferencji i~analizy granicy Pareto. W~rozdziale piątym opisano architekturę i~implementację rozszerzenia przeglądarki Chrome. Rozdział szósty zawiera wnioski końcowe, weryfikację hipotezy i~propozycje dalszego rozwoju.

\chapter{Podstawy teoretyczne i przegląd literatury}

\section{Zjawisko dezinformacji: definicje, typologia i charakterystyka fake newsów}

Dezinformacja jest jednym z~kluczowych wyzwań współczesnych systemów informacyjnych. W~literaturze definiowana jest jako celowe tworzenie i~rozpowszechnianie fałszywych lub wprowadzających w~błąd treści w~celu wywierania wpływu na opinię publiczną albo osiągania korzyści politycznych, ekonomicznych bądź społecznych~\cite{Shu2017}.

W~badaniach tego zjawiska zasadnicze znaczenie ma rozróżnienie między dezinformacją (\textit{disinformation}) a~misinformacją (\textit{misinformation}). Dezinformacja jest działaniem intencjonalnym, natomiast misinformacja obejmuje niezamierzone rozpowszechnianie nieprawdziwych treści~\cite{Hu2025, Yi2025}. Podział ten determinuje charakter danych wykorzystywanych w~systemach automatycznej detekcji oraz przyjmowane założenia modelowe.

Fake newsy stanowią szczególną kategorię dezinformacji: są to treści imitujące format przekazów dziennikarskich, lecz pozbawione faktycznego oparcia w~rzeczywistości. Stosowanie dziennikarskich konwencji -- takich jak nagłówek, lead i~cytat -- utrudnia ich identyfikację przez czytelników~\cite{Hu2025}. W~mediach społecznościowych mechanizmy udostępniania i~algorytmiczne rekomendacje treści dodatkowo sprzyjają wirusowemu rozprzestrzenianiu fałszywych narracji.

Typologia fake newsów wyróżnia treści całkowicie fałszywe, informacje zmanipulowane kontekstowo (np.\ przez zmianę podpisu pod materiałem graficznym), treści częściowo prawdziwe ze~selektywnie przedstawionymi faktami, a~także publikacje satyryczne, które -- wyrwane z~kontekstu -- mogą być mylnie odbierane jako rzetelne informacje~\cite{Hu2025}. Szczególnym przypadkiem są tzw.\ \textit{toxic fake newsy}, łączące dezinformację z~treściami emocjonalnie obciążającymi lub szkodliwymi społecznie; zjawisko to zyskało na znaczeniu w~trakcie pandemii COVID-19, gdy fałszywe informacje zdrowotne mogły bezpośrednio kształtować zachowania społeczeństwa~\cite{Wani2024}.

\section{Stan wiedzy w zakresie detekcji fake newsów}

Rozwój metod automatycznej detekcji fake newsów przebiega równolegle z~postępem w~dziedzinie przetwarzania języka naturalnego i~uczenia maszynowego. Wczesne podejścia opierały się na klasycznych algorytmach -- Naive Bayes, SVM, Random Forest -- które wykorzystywały statystyczne cechy tekstu, m.in.\ częstości n-gramów i~reprezentacje TF-IDF~\cite{Jouhar2024}.

Przełomem okazało się wprowadzenie gęstych reprezentacji wektorowych słów (\textit{word embeddings}), takich jak Word2Vec i~GloVe. Odwzorowując słowa jako punkty w~przestrzeni semantycznej, umożliwiły one uwzględnienie relacji znaczeniowych między wyrazami. Al-Tarawneh i~in.~\cite{AlTarawneh2024} wykazali, że zastosowanie embeddingów poprawia skuteczność modeli głębokiego uczenia, w~szczególności sieci konwolucyjnych (CNN) i~rekurencyjnych (RNN).

Współcześnie dominującym podejściem w~detekcji fake newsów są modele oparte na architekturze Transformer. Centralną rolę odgrywa BERT~\cite{Devlin2019BERTPO}, który dzięki dwukierunkowemu mechanizmowi uwagi umożliwia głębsze modelowanie kontekstu semantycznego. Na tej podstawie opracowano wyspecjalizowane architektury: FakeBERT łączy reprezentacje BERT z~warstwami konwolucyjnymi w~celu równoczesnego uchwycenia kontekstu globalnego i~lokalnych wzorców językowych~\cite{Kaliyar2021}, a~Enhanced BERT wprowadza dodatkowe modyfikacje optymalizujące klasyfikację treści informacyjnych~\cite{Aljawarneh2024}.

Najnowsze kierunki badań obejmują podejścia hybrydowe integrujące modele głębokiego uczenia z~dużymi modelami językowymi (LLM). Teo i~in.~\cite{Teo2024} wskazują, że połączenie LLM z~klasycznymi metodami uczenia maszynowego pozwala czerpać zarówno ze~stabilności modeli tradycyjnych, jak i~ze~zdolności LLM do analizy kontekstowej. Rozwijane są również podejścia agentowe, w~których model językowy nie tylko klasyfikuje tekst, lecz aktywnie weryfikuje informacje na podstawie zewnętrznych źródeł wiedzy~\cite{Li2024}.

\section{Identyfikacja luk badawczych i ograniczeń istniejących metod}

Pomimo intensywnego rozwoju metod detekcji fake newsów literatura wskazuje na szereg istotnych ograniczeń. Pierwszym z~nich jest zjawisko dryfu koncepcji (\textit{concept drift}): modele trenowane na historycznych danych tracą skuteczność wobec nowych form dezinformacji, których cechy językowe i~retoryczne ewoluują~\cite{Hu2025}.

Drugi problem dotyczy interpretowalności modeli głębokiego uczenia. Mimo wysokiej skuteczności klasyfikacyjnej, architektury takie jak BERT czy CNN nie dostarczają transparentnych wyjaśnień podejmowanych decyzji, co ogranicza ich zastosowanie w~domenach wymagających wysokiego poziomu zaufania~\cite{Alshuwaier2025}.

Osobnym wyzwaniem jest rosnąca jakość treści generowanych przez LLM, których styl i~spójność mogą być nie do odróżnienia od rzetelnych artykułów. Yi i~Huang~\cite{Yi2025} zwracają uwagę, że proliferacja syntetycznej dezinformacji wymaga opracowania nowych metod detekcji wykraczających poza analizę cech stylistycznych.

Wreszcie, jakość i~reprezentatywność dostępnych zbiorów danych często nie odzwierciedlają rzeczywistego rozkładu zjawiska -- modele trenowane na wąsko zdefiniowanych korpusach mogą wykazywać ograniczoną zdolność generalizacji poza dystrybucję treningową~\cite{Shu2017}. Niezbędne są badania nad metodami odpornymi na zmienność danych oraz integrującymi różnorodne źródła informacji~\cite{Teo2024}.

\chapter{Metodologia}

\section{Charakterystyka wykorzystanych zbiorów danych}

Badania przeprowadzono z~wykorzystaniem trzech powszechnie stosowanych zbiorów referencyjnych: ISOT, LIAR oraz WELFake. Poniżej przedstawiono szczegółową charakterystykę każdego z~nich, obejmującą pochodzenie, strukturę, podział ilościowy oraz sposób etykietowania.

\subsection{Zbiór danych ISOT}

\subsubsection{Pochodzenie i charakterystyka danych.} 
Zbiór ISOT (\textit{Information Security and Object Technology}) opracowano na Uniwersytecie Wiktorii. Artykuły prawdziwe pozyskano z~serwisu agencji Reuters, a~fałszywe -- ze~stron zidentyfikowanych przez organizację PolitiFact jako źródła dezinformacji. Większość tekstów odnosi się do wydarzeń z~lat 2016--2017, w~szczególności do amerykańskich wyborów prezydenckich.

\subsubsection{Struktura i atrybuty zbioru.} 
Zbiór udostępniony jest w~postaci dwóch plików CSV (dla artykułów prawdziwych i~fałszywych). Każdy rekord zawiera tytuł artykułu (\textit{title}), pełną treść (\textit{text}), kategorię tematyczną (\textit{subject}) oraz datę publikacji (\textit{date}).

\subsubsection{Podział ilościowy i tematyczny.} 
Zbiór jest stosunkowo zbalansowany, z~niewielką przewagą artykułów fałszywych. Dominującą tematyką jest polityka i~wydarzenia ze~świata, co przedstawiono szczegółowo w~Tabeli~\ref{tab:isot_stat}.

\begin{table}[htbp]
\centering
\caption{Podział ilościowy i tematyczny artykułów w zbiorze ISOT}
\vspace{0.2cm}
\begin{tabularx}{\textwidth}{l c X}
\toprule
\textbf{Typ wiadomości} & \textbf{Liczba artykułów} & \textbf{Tematyka} \\
\midrule
Prawdziwe (Real-News) & 21 417 & World-News, Politics-News \\
Fałszywe (Fake-News) & 23 481 & Politics, News, Left-news, Government-News, US News, Middle-east \\
\bottomrule
\end{tabularx}
\label{tab:isot_stat}
\end{table}

\subsubsection{Sposób etykietowania.} 
Problem klasyfikacji w przypadku zbioru ISOT ma charakter ściśle binarny. Etykiety zostały przypisane na podstawie źródła pochodzenia artykułu: teksty z agencji Reuters oznaczono jako prawdziwe (klasa~0), a artykuły ze stron zidentyfikowanych przez PolitiFact jako szerzące dezinformację -- jako fałszywe (klasa~1).

\subsection{Zbiór danych LIAR}

\subsubsection{Pochodzenie i charakterystyka danych.} 
Zbiór LIAR, wprowadzony przez Wanga~\cite{wang2017liar}, należy do najpowszechniej stosowanych benchmarków detekcji dezinformacji. W~odróżnieniu od ISOT, nie zawiera pełnych artykułów, lecz krótkie wypowiedzi (ang.\ \textit{statements}) polityków, publicystów i~osób publicznych, pozyskane z~serwisu fact-checkingowego PolitiFact.com.

\subsubsection{Struktura i atrybuty zbioru.} 
Zbiór wyróżnia się bogatą warstwą metadanych: oprócz treści wypowiedzi zawiera informacje o~mówcy, przynależności partyjnej, kontekście wypowiedzi (np.\ debata telewizyjna, wpis na Twitterze) oraz historii weryfikacji (\textit{credit history}), co umożliwia analizę wykraczającą poza sam tekst.

\subsubsection{Podział ilościowy i tematyczny.} 
Zbiór liczy łącznie 12~817 wypowiedzi, obejmujących tematykę gospodarczą, zdrowotną, podatkową i~społeczną w~Stanach Zjednoczonych. Strukturę binarną przedstawiono w~Tabeli~\ref{tab:liar_stat}.

\begin{table}[htbp]
\centering
\caption{Podział ilościowy i tematyczny wypowiedzi w zbiorze LIAR}
\vspace{0.2cm}
\begin{tabularx}{\textwidth}{l c X}
\toprule
\textbf{Typ wiadomości} & \textbf{Liczba tekstów} & \textbf{Tematyka} \\
\midrule
Prawdziwe & 7 152 & Amerykańska polityka, Gospodarka, Zdrowie, Społeczeństwo \\
Fałszywe & 5 665 & Amerykańska polityka, Gospodarka, Zdrowie, Społeczeństwo \\
\bottomrule
\end{tabularx}
\label{tab:liar_stat}
\end{table}

\subsubsection{Sposób etykietowania.}
\label{sec:liar_label}
Oryginalnie zbiór LIAR wyróżnia sześć stopni prawdziwości (tzw.\ \textit{truth-o-meter}): \textit{pants-on-fire, false, barely-true, half-true, mostly-true} oraz \textit{true}. Na potrzeby niniejszej pracy zbiór poddano binaryzacji: etykiety wskazujące na nieprawdziwość wypowiedzi (\textit{pants-on-fire, false, barely-true}) zmapowano do klasy~1, a etykiety oznaczające prawdziwość (\textit{half-true, mostly-true, true}) -- do klasy~0.

\subsection{Zbiór danych WELFake}

\subsubsection{Pochodzenie i charakterystyka danych.} 
Zbiór WELFake (\textit{Word Embedding over Linguistic Features for Fake News Detection}) powstał poprzez agregację czterech zbiorów: Kaggle, McIntire, Reuters oraz BuzzFeed Political~\cite{verma2021welfake}. Celem tej konsolidacji było uzyskanie dużego, zróżnicowanego korpusu redukującego ryzyko nadmiernego dopasowania (\textit{overfitting}), na które narażone są mniejsze zbiory jednoźródłowe.

\subsubsection{Struktura i atrybuty zbioru.} 
Ze względu na wieloźródłowość dane zostały zunifikowane do wspólnego schematu: identyfikator, tytuł artykułu (\textit{title}), pełna treść (\textit{text}) i~etykieta (\textit{label}).

\subsubsection{Podział ilościowy i tematyczny.} 
Zbiór liczy ponad 72~tysiące artykułów i~cechuje się bardzo dobrym wyważeniem klas. Tematycznie jest różnorodny -- obejmuje wiadomości krajowe i~światowe z~wielu dziedzin. Szczegóły przedstawiono w~Tabeli~\ref{tab:welfake_stat}.

\begin{table}[htbp]
\centering
\caption{Podział ilościowy i tematyczny artykułów w zbiorze WELFake}
\vspace{0.2cm}
\begin{tabularx}{\textwidth}{l c X}
\toprule
\textbf{Typ wiadomości} & \textbf{Liczba artykułów} & \textbf{Tematyka} \\
\midrule
Prawdziwe (Real-News) & 35 028 & Polityka, Wiadomości krajowe i światowe (agregacja) \\
Fałszywe (Fake-News) & 37 106 & Polityka, Wiadomości krajowe i światowe (agregacja) \\
\bottomrule
\end{tabularx}
\label{tab:welfake_stat}
\end{table}

\subsubsection{Sposób etykietowania.} 
Etykietowanie ma charakter binarny, wynikający ze~struktury połączonych zbiorów. Klasa~0 oznacza artykuły fałszywe, klasa~1 -- prawdziwe.

\section{Proces przygotowania danych}
Surowe dane tekstowe rzadko nadają się do bezpośredniego wykorzystania w procesie uczenia maszynowego. Zakłócenia takie jak błędy formatowania, nadmiarowe spacje, adresy URL czy zbędne znaki interpunkcyjne mogą negatywnie wpływać na jakość tworzonych klasyfikatorów, a odpowiedni dobór technik wstępnego przetwarzania (ang.\ \textit{data preprocessing}) ma krytyczne znaczenie dla ich ostatecznej skuteczności~\cite{AIZAWA200345}. Jednocześnie najnowsze badania wskazują, że wpływ tych technik różni się diametralnie w zależności od rodzaju modelu: klasyczne algorytmy i nowoczesne architektury Transformer reagują na preprocessing w odmienny sposób~\cite{SIINO2024102342}. Z tego względu etap przygotowania danych podzielono na trzy wyraźne kroki chronologiczne: wstępne czyszczenie i unifikację etykiet, rygorystyczny podział na podzbiory ewaluacyjne, a na końcu -- dedykowaną normalizację tekstu dostosowaną do wymagań konkretnego algorytmu.

\subsection{Czyszczenie specyficzne i obsługa etykiet}

Pierwszym krokiem, wykonywanym przed podziałem na podzbiory, było ujednolicenie struktury zbiorów i~usunięcie rekordów z~brakującymi wartościami w~kolumnie tekstowej. Ze~względu na różnice w~formatach wejściowych każdy zbiór wymagał dedykowanych operacji:
\begin{itemize}
    \item \textbf{Zbiór LIAR:} Sześć oryginalnych etykiet zmapowano do postaci binarnej zgodnie z~opisem w~sekcji~\ref{sec:liar_label}.
    \item \textbf{Zbiór ISOT:} Artykuły prawdziwe pochodziły wyłącznie od agencji Reuters, co groziło wyciekiem danych (\textit{data leakage}) -- model mógłby nauczyć się klasyfikować teksty jako prawdziwe wyłącznie na podstawie obecności słowa ,,Reuters''. Wdrożono moduł czyszczący, który identyfikuje i~usuwa tę nazwę własną oraz zbędne znaczniki multimedialne (np.\ \texttt{[video]}).
    \item \textbf{Zbiór WELFake:} Dane zweryfikowano pod kątem spójności strukturalnej i~usunięto puste rekordy. Etykiety zachowano w~oryginalnym formacie (0~=~fałszywe, 1~=~prawdziwe).
\end{itemize}

Należy zaznaczyć, że zastosowane mapowania etykiet są wzajemnie niespójne: ISOT i~LIAR używają konwencji \textit{Fake}~=~1, \textit{Real}~=~0, podczas gdy WELFake zachowuje oryginalną konwencję \textit{Fake}~=~0, \textit{Real}~=~1. Metryki ewaluacyjne obliczane na zbiorach testowych (F1, precyzja, czułość) są poprawne dla każdego zbioru z~osobna, ponieważ sklearn porównuje predykcje z~etykietami w~ramach tej samej konwencji. Niespójność ta może natomiast wpływać na kierunek predykcji modeli ISOT i~LIAR udostępnionych przez API (zob.\ sekcja~\ref{sec:impl_router}) i~stanowi obszar wymagający korekty w~kolejnej iteracji projektu.

\subsection{Podział na podzbiory uczący, walidacyjny i testowy}

Po ujednoliceniu danych podzielono je na trzy rozłączne podzbiory, eliminując ryzyko wycieku informacji między etapami przetwarzania:
\begin{itemize}
    \item \textbf{Zbiór uczący (\textit{Train})} -- podstawa do treningu algorytmów i~aktualizacji wag.
    \item \textbf{Zbiór walidacyjny (\textit{Validation})} -- służy do strojenia hiperparametrów i~wczesnego zatrzymywania treningu (\textit{early stopping}), ograniczając ryzyko przeuczenia.
    \item \textbf{Zbiór testowy (\textit{Test})} -- zbiór ,,niewidziany'' przez modele podczas treningu i~strojenia, przeznaczony wyłącznie do ostatecznej ewaluacji.
\end{itemize}

Dla zbiorów ISOT i~WELFake dane przetasowano i~podzielono w~proporcjach \textbf{80/10/10}. Zbiór LIAR zachował oryginalny podział przygotowany przez autorów, co umożliwia bezpośrednie porównywanie wyników z~innymi pracami. We wszystkich losowych podziałach zastosowano ziarno losowości (\textit{random state}~=~42) oraz \textbf{próbkowanie warstwowe} (\textit{stratified sampling}), gwarantujące zachowanie proporcji klas we wszystkich trzech podzbiorach.

\subsection{Wstępne przetwarzanie tekstu}

Przed podaniem danych do modeli stosowano jeden z~dwóch dedykowanych preprocesorów, dobieranych na podstawie rodzaju klasyfikatora.

\subsubsection*{Tryb klasyczny (\textit{Classic Mode})}

Klasyczne algorytmy oparte na rzadkich reprezentacjach (TF-IDF, BoW) są wrażliwe na szum informacyjny i~obecność rzadkich tokenów~\cite{AIZAWA200345}. Agresywne czyszczenie tekstu ogranicza wymiarowość przestrzeni cech i~poprawia jakość klasyfikacji. Zastosowany potok normalizacji obejmuje kolejno:
\begin{itemize}
    \item zamianę wszystkich znaków na małe litery (\textit{lowercasing}),
    \item usunięcie adresów URL (ciągów zaczynających się od \texttt{http}, \texttt{https} lub \texttt{www}),
    \item usunięcie znaczników HTML,
    \item eliminację znaków interpunkcyjnych i~cyfr -- pozostawiane są wyłącznie ciągi alfabetyczne,
    \item normalizację białych znaków -- wielokrotne spacje, tabulatory i~znaki nowej linii zastępowane są pojedynczą spacją.
\end{itemize}

\subsubsection*{Tryb dla modeli Transformer (\textit{BERT Mode})}

Modele z~rodziny Transformer opierają się na mechanizmie samouwagi (\textit{self-attention}), dla którego agresywne czyszczenie prowadzi do utraty cennego kontekstu semantycznego i~syntaktycznego~\cite{camachocollados2018role}. Interpunkcja, wielkie litery i~cyfry niosą w~tych architekturach istotną informację (np.\ charakterystyczne dla dezinformacji użycie capslocka czy pytania retoryczne). Współczesne badania potwierdzają, że dla zaawansowanych modeli językowych ograniczenie preprocessingu do minimum przynosi najlepsze rezultaty~\cite{SIINO2024102342}. Dlatego w~trybie BERT zachowano oryginalną kapitalizację, interpunkcję i~liczby, wykonując jedynie \textbf{normalizację białych znaków}.

\section{Metody reprezentacji tekstu}
Algorytmy uczenia maszynowego wymagają numerycznej reprezentacji tekstu. W~niniejszym projekcie zastosowano trzy kategorie metod wektoryzacji: statystyczne, gęste reprezentacje wektorowe oraz kontekstowe embeddingi transformerowe.

\subsection{Podejścia statystyczne}

Podstawową metodą statystyczną jest model worka słów (\textit{Bag of Words}, BoW), który reprezentuje dokument jako wektor zliczeń tokenów, ignorując ich kolejność i~strukturę gramatyczną~\cite{8950616}. Prostota BoW przekłada się na szybkość, ale kosztem rzadkich, wysokowymiarowych wektorów niemożliwych do uchwycenia synonimii ani wieloznaczności.

Bardziej wyrafinowanym rozwinięciem jest metoda TF-IDF (\textit{Term Frequency-Inverse Document Frequency}), która nadaje tokenom wagi na podstawie ich lokalnej częstości w~dokumencie korygowanej przez odwrotność częstości w~całym korpusie~\cite{AIZAWA200345}. Powszechne słowa funkcyjne (spójniki, przyimki) otrzymują niskie wagi, natomiast rzadkie terminy charakteryzujące dany dokument -- wysokie. W~projekcie obie metody zaimplementowano jako klasy \texttt{BowEmbedder} oraz \texttt{TfidfEmbedder}.

\subsection{Reprezentacja wektorowa}
Ograniczenia rzadkich reprezentacji statystycznych skłoniły do opracowania gęstych wektorów słów (\textit{word embeddings}). Reprezentują one słowa jako punkty w~ciągłej przestrzeni wielowymiarowej, w~której podobieństwo wektorów odzwierciedla podobieństwo semantyczne wyrazów.

Pionierskim modelem w~tej dziedzinie jest Word2Vec~\cite{mikolov2013efficient}, który za pomocą płytkich sieci neuronowych (architektury CBOW lub Skip-gram) uczy się reprezentacji na podstawie lokalnego kontekstu wyrazów~\cite{CHURCH_2017}. Alternatywę stanowi GloVe (\textit{Global Vectors for Word Representation})~\cite{pennington2014glove}, opierający się na faktoryzacji globalnej macierzy współwystępowania słów w~całym korpusie. GloVe łączy korzyści globalnej analizy statystycznej z~lokalnym uczeniem kontekstowym. W~projekcie oba podejścia zintegrowano jako klasy \texttt{Word2VecEmbedder} i~\texttt{GloveEmbedder}.

\subsection{Embeddingi kontekstowe}
\label{sec:embeddingi_kontekstowe}
Statyczność tradycyjnych embeddingów -- każde słowo ma dokładnie jedną reprezentację niezależnie od kontekstu -- uniemożliwia rozróżnienie polisemii (np.\ słowo ,,zamek'' jako budowla vs.\ element zapięcia). Problem ten rozwiązują kontekstowe embeddingi generowane przez modele Transformer.

Modele te wyznaczają reprezentację każdego tokenu dynamicznie, na podstawie całej sekwencji wejściowej. W~projekcie wykorzystano cztery modele z~rodziny BERT:
\begin{itemize}
    \item \textbf{BERT} (\texttt{bert-base-uncased}) -- model bazowy uczący się dwukierunkowych reprezentacji tekstu poprzez mechanizm samouwagi (\textit{self-attention})~\cite{Devlin2019BERTPO}.
    \item \textbf{RoBERTa} (\texttt{roberta-base}) -- zoptymalizowana wersja BERT z~rozszerzonym procesem pretreningu na większych zbiorach danych, oferująca lepszą generalizację~\cite{Liu2019RoBERTa}.
    \item \textbf{DistilBERT} (\texttt{distilbert-base-uncased}) -- model skompresowany metodą destylacji wiedzy, zachowujący większość skuteczności BERT przy znacząco mniejszej liczbie parametrów i~krótszym czasie inferencji~\cite{sanh2020distilbert}.
    \item \textbf{FakeBERT} (\texttt{bert-base-uncased}) -- wariant zainspirowany architekturą zaproponowaną przez Kaliyara i~in.~\cite{Kaliyar2021}, zaimplementowany w~niniejszym projekcie jako dostrajany model \textit{bert-base-uncased} z~głową klasyfikacyjną. W~odróżnieniu od oryginalnej publikacji, implementacja nie zawiera dodatkowych warstw konwolucyjnych -- jej celem jest ocena, na ile sama inicjalizacja wagami BERT wpływa na skuteczność detekcji fake newsów.
\end{itemize}
Zastosowanie embeddingów kontekstowych umożliwia uchwycenie subtelnych niuansów semantycznych i~syntaktycznych niedostępnych dla metod statystycznych.


\section{Modele}

W~zadaniach klasyfikacji tekstu, w~tym detekcji fake newsów, stosuje się modele o~bardzo zróżnicowanych architekturach -- od algorytmów statystycznych, przez sieci głębokie, po transformery i~wielkie modele językowe~\cite{Jouhar2024, Alshuwaier2025}. Poniżej opisano trzy rodziny algorytmów zastosowanych w~badaniach.

\subsection{Modele klasycznego uczenia maszynowego}

Klasyczne algorytmy uczenia maszynowego wymagają uprzedniej transformacji tekstu do postaci numerycznej -- najczęściej przy użyciu BoW lub TF-IDF. Regresja logistyczna~\cite{Dreiseitl2002} i~naiwny klasyfikator Bayesa~\cite{Murphy2006NB} cenione są za prostotę, szybkość i~interpretowalność. Maszyny wektorów nośnych (SVM) wyznaczają optymalną hiperpłaszczyznę separującą klasy w~przestrzeni cech~\cite{Abdullah2021SVM, Jakkula2006SVM}, a~$k$-NN klasyfikuje próbki na podstawie podobieństwa do najbliższych sąsiadów~\cite{10.1007/978-3-540-39964-3_62}.

Wysoką skuteczność osiągają także modele drzewiaste~\cite{barry2013decision}, a~zwłaszcza metody zespołowe: lasy losowe (\textit{Random Forests})~\cite{Breiman2001} i~XGBoost~\cite{Chen2016XGBoost}. Łącząc wyniki wielu słabszych klasyfikatorów bazowych, ograniczają problem nadmiernego dopasowania i~poprawiają generalizację.

\subsection{Modele uczenia głębokiego}

Konieczność ręcznej inżynierii cech w~klasycznych metodach ML ogranicza ich zdolność do uchwycenia złożonych wzorców językowych. Modele głębokiego uczenia~\cite{Alzubaidi2021} automatycznie uczą się hierarchicznych reprezentacji danych, efektywniej kodując zależności semantyczne i~składniowe~\cite{AlTarawneh2024}.

Konwolucyjne sieci neuronowe (CNN) -- pierwotnie opracowane dla analizy obrazów -- znalazły zastosowanie w~ekstrakcji lokalnych wzorców językowych (n-gramów i~charakterystycznych układów słów)~\cite{Roumeliotis2025}. Dla danych sekwencyjnych kluczową rolę odgrywają sieci rekurencyjne: LSTM (\textit{Long Short-Term Memory})~\cite{7508408} modelują długoterminowe zależności dzięki bramkowym mechanizmom pamięci, a~BiLSTM~\cite{9005997} analizuje sekwencję w~obu kierunkach, uzyskując pełniejszy obraz kontekstu. Alternatywą są sieci GRU (\textit{Gated Recurrent Unit})~\cite{8053243}, których uproszczona struktura pozwala osiągać porównywalną skuteczność przy niższym koszcie obliczeniowym.

\subsection{Transformery}

Architektura Transformer, bazująca na mechanizmie samouwagi~\cite{Teo2024}, stanowi przełom w~przetwarzaniu języka naturalnego. W~odróżnieniu od modeli rekurencyjnych transformery przetwarzają dane równolegle, efektywnie modelując zależności między dowolnymi elementami sekwencji niezależnie od ich odległości.

Centralnym modelem tej rodziny jest BERT (\textit{Bidirectional Encoder Representations from Transformers})~\cite{Devlin2019BERTPO}, który dzięki wstępnemu treningowi na dużych korpusach i~dwukierunkowemu mechanizmowi uwagi znacząco poprawił skuteczność wielu zadań NLP~\cite{Aljawarneh2024}. Jego warianty optymalizują różne aspekty tej architektury: RoBERTa~\cite{Liu2019RoBERTa} wydłuża pretrening na większych zbiorach, uzyskując lepszą generalizację, a~DistilBERT~\cite{sanh2020distilbert} stosuje destylację wiedzy, redukując liczbę parametrów przy zachowaniu wysokiej jakości reprezentacji.

Dla zadania detekcji fake newsów Kaliyar i~in.~\cite{Kaliyar2021} zaproponowali FakeBERT -- architekturę łączącą reprezentacje BERT z~jednowymiarowymi warstwami konwolucyjnymi, zaprojektowaną z~myślą o~analizie wiarygodności treści medialnych. Jak opisano w~sekcji~\ref{sec:embeddingi_kontekstowe}, w~niniejszym projekcie FakeBERT został zaimplementowany jako model fine-tunowany na identycznej ścieżce treningowej co pozostałe warianty BERT.

Najnowsze badania coraz częściej sięgają po wielkie modele językowe (LLM), które potrafią wykonywać złożoną analizę semantyczną, wykrywać manipulacje językowe i~wspierać weryfikację informacji przez analizę zewnętrznych źródeł wiedzy~\cite{Li2024, Yi2025}.

\section{Projekt architektury modeli eksperymentalnych}

W~celu powtarzalnej i~skalowalnej ewaluacji algorytmów detekcji fake newsów zaprojektowano zautomatyzowany potok eksperymentalny zaimplementowany w~skrypcie \texttt{research/train\_models.py}. System iteracyjnie przeprowadza eksperymenty dla wszystkich kombinacji modeli, reprezentacji tekstu i~zbiorów danych (ISOT, LIAR, WELFake)~\cite{kaggle_isot, kaggle_liar, kaggle_welfake}, opierając się na trzech filarach: deklaratywnym zarządzaniu konfiguracją, monitorowaniu eksperymentów i~automatycznym strojeniu hiperparametrów.

\subsubsection*{Zarządzanie konfiguracją i monitorowanie -- Hydra i MLflow}

Konfigurację eksperymentów zarządza framework Hydra~\cite{Yadan2019Hydra}, który umożliwia hierarchiczną organizację parametrów w~plikach YAML i~ich nadpisywanie z~wiersza poleceń bez ingerencji w~kod źródłowy. Przebieg eksperymentów, metryki ewaluacyjne i~wytrenowane modele logowane są w~lokalnej instancji MLflow na bazie SQLite (\texttt{mlflow.db}).

\subsubsection*{Automatyczne strojenie hiperparametrów -- Optuna}

Potok rozbudowano o~moduł strojenia hiperparametrów oparty na bibliotece Optuna~\cite{akiba2019optuna}. Architektura \textit{define-by-run} pozwala na dynamiczne konstruowanie przestrzeni poszukiwań. Jako algorytm próbkowania zastosowano TPE (\textit{Tree-structured Parzen Estimator}), a~do wczesnego zatrzymywania nieobiecujących prób -- \textit{median pruner}.

\subsubsection*{Przestrzenie poszukiwań i optymalizowane hiperparametry}

Dla każdej rodziny algorytmów zdefiniowano odrębne przestrzenie poszukiwań:

\begin{itemize}
    \item \textbf{Klasyczne algorytmy ML (SVM, RF, XGBoost i~in.):}
    10 prób na eksperyment. Dla SVM: parametr regularyzacji $C \in [0{,}1;\,100]$ (skala logarytmiczna) i~maksymalna liczba iteracji. Dla modeli drzewiastych (RF, XGB): liczba estymatorów $\in [50;\,300]$, maksymalna głębokość drzew $\in [3;\,12]$ oraz współczynnik uczenia.

    \item \textbf{Sieci głębokie (LSTM, GRU, BiLSTM, CNN):}
    8 prób na architekturę. Optymalizowane parametry: rozmiar warstwy ukrytej (64, 128 lub 256 neuronów), liczba warstw (1 lub 2), \textit{batch size} (32 lub 64), \textit{dropout} $\in [0{,}1;\,0{,}5]$ i~\textit{learning rate} $\in [10^{-4};\,10^{-2}]$ (skala logarytmiczna).

    \item \textbf{Transformery (BERT, RoBERTa, DistilBERT, FakeBERT):}
    10 prób ze względu na złożoność obliczeniową. Przestrzeń poszukiwań obejmuje \textit{learning rate} $\in [10^{-5};\,5{\cdot}10^{-5}]$, \textit{batch size} (16 lub 32), maksymalną długość sekwencji (128, 256 lub 512 tokenów), \textit{weight decay} $\in [0;\,0{,}1]$, \textit{warmup ratio} $\in [0;\,0{,}2]$, typ schedulera (\textit{linear} lub \textit{cosine}) oraz liczbę zamrożonych warstw kodera (\textit{freeze\_layers} $\in [0;\,8]$), co pozwala zapobiec zapominaniu wiedzy z~pretreningu.
\end{itemize}


\subsection{Mapa kompatybilności reprezentacji tekstu}
Różne rodziny modeli mają odmienne wymagania co do formatu danych wejściowych. Zdefiniowano macierz kompatybilności (\texttt{COMPATIBILITY\_MAP}) określającą, które metody wektoryzacji można łączyć z~konkretnymi klasyfikatorami:
\begin{itemize}
    \item \textbf{Modele klasyczne} (SVM, LR, RF, DT, XGBoost) ewaluowane są z~reprezentacjami TF-IDF i~BoW oraz gęstymi wektorami Word2Vec i~GloVe. Wyjątkiem jest wielomianowy naiwny Bayes (MNB), który wymaga nieujemnych wartości wejściowych, a~zatem współpracuje wyłącznie z~TF-IDF i~BoW.
    \item \textbf{Sieci głębokie} (LSTM, GRU, BiLSTM, CNN) inicjalizowane są wstępnie wyuczonymi macierzami osadzeń GloVe i~Word2Vec, z~których konstruowane są wewnętrzne warstwy embeddingów.
    \item \textbf{Transformery} (BERT, RoBERTa, FakeBERT, DistilBERT) korzystają z~dedykowanych tokenizatorów i~wag inicjalizacyjnych powiązanych z~danym checkpointem (np.\ \texttt{bert-base-uncased}, \texttt{roberta-base}).
\end{itemize}

\subsection{Proces optymalizacji hiperparametrów}
Automatyczne przeszukiwanie przestrzeni hiperparametrów realizuje biblioteka Optuna~\cite{akiba2019optuna}. Plik konfiguracyjny \texttt{config.yaml} definiuje algorytm próbkowania TPE (\textit{Tree-structured Parzen Estimator}) oraz strategię wczesnego zatrzymywania nieefektywnych prób (\textit{median pruner}). Liczba prób jest dostosowana do zasobochłonności modeli: 10 dla transformerów i~klasycznych algorytmów ML oraz 8 dla sieci głębokich. Szczegółowe przestrzenie poszukiwań dla każdej rodziny architektur zostały opisane w~poprzedniej sekcji. Tak skonfigurowany potok umożliwia rzetelne, powtarzalne porównanie dziesiątek kombinacji model--reprezentacja w~ujednoliconych warunkach eksperymentalnych.

\section{Dobór metryk ewaluacji}

Ewaluacja modeli opiera się na zadaniu klasyfikacji binarnej: model przypisuje tekst do klasy \textbf{Real} (prawdziwy) lub \textbf{Fake} (fałszywy)~\cite{zhou2020survey, Shu2017}. Zastosowano cztery standardowe metryki~\cite{tharwat2020classification}: dokładność (\textit{accuracy}), precyzję (\textit{precision}), czułość (\textit{recall}) i~miarę F1 -- stanowiące uznany standard w~detekcji dezinformacji~\cite{zhou2020survey, Hu2025}.

Podstawą do wyliczenia tych wskaźników jest \textbf{macierz pomyłek} (\textit{confusion matrix})~\cite{tharwat2020classification}. W~kontekście detekcji fake newsów jej elementy interpretuje się następująco:
\begin{itemize}
    \item \textbf{TP (\textit{True Positives})} -- artykuły fałszywe poprawnie sklasyfikowane jako \textit{Fake}.
    \item \textbf{TN (\textit{True Negatives})} -- artykuły prawdziwe poprawnie sklasyfikowane jako \textit{Real}.
    \item \textbf{FP (\textit{False Positives})} -- błąd pierwszego rodzaju: artykuł prawdziwy błędnie oznaczony jako \textit{Fake}.
    \item \textbf{FN (\textit{False Negatives})} -- błąd drugiego rodzaju: artykuł fałszywy niesłusznie uznany za \textit{Real}.
\end{itemize}

Na podstawie powyższych elementów wyznaczono cztery docelowe miary~\cite{tharwat2020classification}:

\textbf{1. Dokładność (\textit{accuracy})} -- odsetek wszystkich poprawnych predykcji (TP i~TN) względem łącznej liczby przypadków. W~zbiorach niezbalansowanych może być zwodnicza, faworyzując klasę dominującą~\cite{powers2011evaluation}:
\[ Accuracy = \frac{TP + TN}{TP + TN + FP + FN} \]

\textbf{2. Precyzja (\textit{precision})} -- jaki odsetek przypadków oznaczonych przez model jako \textit{Fake} był rzeczywiście fałszywy. Istotna w~systemach, gdzie koszt fałszywego alarmu (cenzurowanie rzetelnego artykułu) jest wysoki:
\[ Precision = \frac{TP}{TP + FP} \]

\textbf{3. Czułość (\textit{recall})} -- jaka część wszystkich faktycznie fałszywych artykułów w~zbiorze testowym została poprawnie wykryta. Kluczowa, gdy priorytetem jest minimalizacja zasięgu dezinformacji:
\[ Recall = \frac{TP}{TP + FN} \]

\textbf{4. Miara F1 (\textit{F1-score})} -- średnia harmoniczna precyzji i~czułości. Ponieważ obie miary są ze sobą odwrotnie zależne~\cite{powers2011evaluation}, F1 promuje modele utrzymujące ich zrównoważony, wysoki poziom~\cite{tharwat2020classification}:
\[ F1 = 2 \cdot \frac{Precision \cdot Recall}{Precision + Recall} \]

Miara F1 stanowi główne kryterium rankingu w~analizie eksperymentalnej~\cite{zhou2020survey, Hu2025}; remisy rozstrzyga dokładność.


\chapter{Eksperymenty i analiza porównawcza wyników}

\section{Wyniki klasycznych modeli uczenia maszynowego}
\label{sec:wyniki_klasyczne}

W~rozdziale tym zestawiono wyniki ośmiu klasycznych algorytmów uczenia
maszynowego: regresji logistycznej (\textit{lr}), maszyny wektorów
nośnych (\textit{svm}), drzewa decyzyjnego (\textit{dt}), lasu
losowego (\textit{rf}), XGBoost (\textit{xgb}), naiwnego klasyfikatora
Bayesa w~wariancie Gaussa (\textit{nb}) i~wielomianowym (\textit{mnb})
oraz $k$-NN (\textit{knn}). Każdy klasyfikator przetestowano we wszystkich
dopuszczalnych kombinacjach z~reprezentacjami TF-IDF, BoW, Word2Vec i~GloVe.
Kombinacje niedozwolone (np.\ MNB z~gęstymi wektorami) były odrzucane
przez fabrykę modeli, stąd liczba eksperymentów w~poszczególnych wierszach
tabel nie jest jednakowa. Hiperparametry dobrano za pomocą Optuna
(próbkowanie TPE), a~ostateczne metryki wyznaczono na zbiorze testowym
(10\,\% każdego ze zbiorów: ISOT, LIAR i~WELFake).

\subsection{Wyniki na zbiorach testowych}
\label{sec:wyniki_klasyczne_test}

Tabela~\ref{tab:klas_isot} zawiera podsumowanie najważniejszych
wyników klasycznych modeli na zbiorze ISOT. Wartości w~tabeli zostały
posortowane malejąco względem miary $F_1$.

\begin{table}[ht]
\centering
\caption{Wyniki klasycznych modeli uczenia maszynowego na zbiorze
testowym ISOT (najlepsze osiem kombinacji wg miary $F_1$).}
\label{tab:klas_isot}
\begin{tabular}{llrrrrrr}
\toprule
Model & Embedding & Acc. & Prec. & Recall & $F_1$ & $t_{tr}$ [s] & $t_{inf}$ [s] \\
\midrule
xgb & tfidf    & 0,9977 & 0,9977 & 0,9977 & 0,9977 & 150,05 & 0,072 \\
xgb & bow      & 0,9974 & 0,9975 & 0,9974 & 0,9974 & 12,14  & 0,063 \\
lr  & tfidf    & 0,9921 & 0,9921 & 0,9921 & 0,9921 & 1,77   & 0,040 \\
svm & tfidf    & 0,9905 & 0,9905 & 0,9905 & 0,9905 & 29,39  & 0,048 \\
svm & word2vec & 0,9836 & 0,9836 & 0,9836 & 0,9836 & 6,60   & 0,023 \\
lr  & word2vec & 0,9813 & 0,9813 & 0,9813 & 0,9813 & 3,30   & 0,016 \\
rf  & tfidf    & 0,9739 & 0,9743 & 0,9739 & 0,9739 & 27,67  & 0,607 \\
mnb & tfidf    & 0,9619 & 0,9619 & 0,9619 & 0,9619 & 0,12   & 0,034 \\
\bottomrule
\end{tabular}
\end{table}

Na zbiorze ISOT klasyczne modele osiągnęły wyjątkowo wysoką skuteczność.
Najlepszy wynik (\(F_1 = 0{,}9977\)) uzyskał XGBoost z~TF-IDF, przy czym
różnica między pierwszą a~ósmą pozycją rankingu wynosi zaledwie 4 punkty
procentowe. Representacje rzadkie (TF-IDF, BoW) wyraźnie dominują nad
gęstymi wektorami uśrednionymi (Word2Vec, GloVe) -- cztery pierwsze miejsca
zajmują wyłącznie modele oparte na TF-IDF lub BoW. Najsłabszy klasyfikator,
$k$-NN z~BoW, uzyskał \(F_1 = 0{,}6444\): wysoka wymiarowość i~rzadkość
przestrzeni cech powoduje, że euklidesowa metryka odległości traci
zdolność do dyskryminacji klas.

Tabela~\ref{tab:klas_welfake} zawiera analogiczne wyniki dla korpusu
WELFake.

\begin{table}[ht]
\centering
\caption{Wyniki klasycznych modeli na zbiorze testowym WELFake
(najlepsze osiem kombinacji wg miary $F_1$).}
\label{tab:klas_welfake}
\begin{tabular}{llrrrrrr}
\toprule
Model & Embedding & Acc. & Prec. & Recall & $F_1$ & $t_{tr}$ [s] & $t_{inf}$ [s] \\
\midrule
xgb & bow      & 0,9713 & 0,9715 & 0,9713 & 0,9713 & 56,17  & 0,173 \\
xgb & tfidf    & 0,9711 & 0,9713 & 0,9711 & 0,9711 & 147,67 & 0,045 \\
lr  & tfidf    & 0,9576 & 0,9576 & 0,9576 & 0,9576 & 8,14   & 0,072 \\
svm & tfidf    & 0,9490 & 0,9490 & 0,9490 & 0,9489 & 58,60  & 0,085 \\
dt  & bow      & 0,9362 & 0,9381 & 0,9362 & 0,9364 & 20,21  & 0,021 \\
dt  & tfidf    & 0,9336 & 0,9349 & 0,9336 & 0,9337 & 57,23  & 0,029 \\
svm & word2vec & 0,9194 & 0,9196 & 0,9194 & 0,9195 & 32,06  & 0,065 \\
lr  & word2vec & 0,9165 & 0,9165 & 0,9165 & 0,9165 & 5,55   & 0,016 \\
\bottomrule
\end{tabular}
\end{table}

Na korpusie WELFake wyniki są niższe niż na ISOT, lecz nadal pozwalają
ocenić zadanie jako rozwiązane w~bardzo dobrym stopniu. Najlepszą
konfiguracją okazał się XGBoost z~BoW (\(F_1 = 0{,}9713\)), potwierdzając
efektywność metod drzewiastych w~parze z~rzadkimi reprezentacjami. Średnia
miara $F_1$ wyniosła \(0{,}883\) -- niżej niż na ISOT, co odzwierciedla
większą trudność zadania: WELFake, jako agregat czterech zbiorów,
charakteryzuje się większą stylistyczną różnorodnością tekstów.

Najtrudniejszym zbiorem okazał się LIAR -- zawiera krótkie wypowiedzi
polityków, w~których sześciostopniowa skala wiarygodności została
zbinaryzowana do problemu dwuklasowego (Tabela~\ref{tab:klas_liar}).

\begin{table}[ht]
\centering
\caption{Wyniki klasycznych modeli na zbiorze testowym LIAR
(najlepsze osiem kombinacji wg miary $F_1$).}
\label{tab:klas_liar}
\begin{tabular}{llrrrrrr}
\toprule
Model & Embedding & Acc. & Prec. & Recall & $F_1$ & $t_{tr}$ [s] & $t_{inf}$ [s] \\
\midrule
mnb & bow      & 0,6335 & 0,6321 & 0,6335 & 0,6326 & 0,01  & 0,008 \\
xgb & bow      & 0,6280 & 0,6315 & 0,6280 & 0,6291 & 3,81  & 0,025 \\
mnb & tfidf    & 0,6312 & 0,6274 & 0,6312 & 0,6232 & 0,01  & 0,013 \\
lr  & glove    & 0,6202 & 0,6280 & 0,6202 & 0,6215 & 1,47  & 0,015 \\
svm & glove    & 0,6186 & 0,6262 & 0,6186 & 0,6199 & 3,32  & 0,020 \\
rf  & tfidf    & 0,6108 & 0,6243 & 0,6108 & 0,6115 & 1,11  & 0,067 \\
nb  & glove    & 0,6092 & 0,6144 & 0,6092 & 0,6105 & 0,03  & 0,010 \\
xgb & tfidf    & 0,6085 & 0,6105 & 0,6085 & 0,6092 & 11,84 & 0,026 \\
\bottomrule
\end{tabular}
\end{table}

Żadna z~klasycznych metod na zbiorze LIAR nie przekroczyła \(F_1 = 0{,}64\).
Najlepszy wynik (\(F_1 = 0{,}6326\)) uzyskał wielomianowy Bayes z~BoW,
co jest typowe dla bardzo krótkich tekstów z~ograniczonym słownikiem.
Wynik bliski poziomowi losowemu (baseline większościowy \(\approx 0{,}56\))
dowodzi, że proste cechy n-gramowe nie wystarczają do uchwycenia subtelności
językowych charakteryzujących fałszywe wypowiedzi polityczne.

\subsection{Analiza czasów uczenia i~wnioskowania}
\label{sec:wyniki_klasyczne_czas}

Tabela~\ref{tab:klas_czasy} zestawia średnie czasy uczenia i~wnioskowania
dla rodziny klasycznej w~podziale na zbiory danych.

\begin{table}[ht]
\centering
\caption{Średnie czasy uczenia i~wnioskowania klasycznych modeli
w~zależności od zbioru danych (w~sekundach).}
\label{tab:klas_czasy}
\begin{tabular}{lrrrr}
\toprule
Zbiór & $\overline{t_{tr}}$ & $\max{t_{tr}}$ & $\overline{t_{inf}}$ & $\max{t_{inf}}$ \\
\midrule
ISOT    & 18,02  & 150,05 & 0,755 & 11,02  \\
LIAR    & 3,21   & 16,38  & 0,044 & 0,234  \\
WELFake & 33,29  & 147,67 & 2,504 & 37,16  \\
\bottomrule
\end{tabular}
\end{table}

Czas uczenia mieści się w~granicach kilkudziesięciu sekund nawet dla
największego zbioru. Najszybciej trenują naiwny Bayes i~drzewa decyzyjne
(ułamki sekundy), najwolniej -- XGBoost. Czas wnioskowania na pełnym
zbiorze testowym zazwyczaj nie przekracza 0{,}1~s, z~wyjątkiem $k$-NN,
którego złożoność liniowo rośnie z~liczbą próbek (do \(37{,}16\)~s dla
WELFake). Potwierdza to niski koszt eksploatacyjny klasycznych modeli
w~porównaniu z~sieciami głębokimi i~transformerami
(por.\ rozdz.~\ref{sec:wyniki_transformery}).

\subsection{Podsumowanie}
\label{sec:wyniki_klasyczne_podsum}

Klasyczne algorytmy uczenia maszynowego stanowią silny i~szybki punkt
odniesienia. Na ISOT najlepsza konfiguracja (XGBoost + TF-IDF)
osiąga niemal idealną skuteczność (\(F_1 = 0{,}9977\)), a~na WELFake
wszystkie czołowe modele przekraczają próg \(0{,}97\). Przy najkrótszych
czasach uczenia i~wnioskowania w~całym badaniu, modele klasyczne
są naturalnymi kandydatami do zastosowań wymagających niskich opóźnień
(np.\ rozszerzenia przeglądarki). Słabym punktem tej rodziny pozostaje
LIAR -- najwyższy osiągnięty wynik to \(F_1 = 0{,}6326\). Zdolność
klasycznych modeli do generalizacji na dane spoza dystrybucji
omówiono w~sekcji~\ref{sec:wyniki_www}.

\section{Wyniki sieci głębokich}
\label{sec:wyniki_glebokie}

W~rozdziale tym omówiono wyniki czterech architektur: LSTM, GRU,
BiLSTM i~jednowymiarowej CNN, zaimplementowanych w~PyTorch.
Hiperparametry (rozmiar warstwy ukrytej, dropout, learning rate,
rozmiar batcha) dobierano przez Optuna. W~odróżnieniu od klasycznych
modeli, sieci głębokie inicjalizują warstwy embeddingów przy użyciu
wstępnie wyuczonych wektorów: Word2Vec (Google News, 300d)
i~GloVe (6B, 100d).

\subsection{Wyniki na zbiorach testowych}
\label{sec:wyniki_glebokie_test}

Tabela~\ref{tab:dl_isot} zawiera podsumowanie wyników sieci głębokich
na zbiorze ISOT.

\begin{table}[ht]
\centering
\caption{Wyniki sieci głębokich na zbiorze testowym ISOT.}
\label{tab:dl_isot}
\begin{tabular}{llrrrrrr}
\toprule
Model & Embedding & Acc. & Prec. & Recall & $F_1$ & $t_{tr}$ [s] & $t_{inf}$ [s] \\
\midrule
gru    & word2vec & 0,9946 & 0,9946 & 0,9946 & 0,9946 & 67,06 & 0,491 \\
bilstm & word2vec & 0,9941 & 0,9941 & 0,9941 & 0,9941 & 71,13 & 0,388 \\
cnn    & word2vec & 0,9916 & 0,9916 & 0,9916 & 0,9916 & 83,04 & 0,425 \\
gru    & glove    & 0,9913 & 0,9913 & 0,9913 & 0,9913 & 55,28 & 0,675 \\
cnn    & glove    & 0,9880 & 0,9882 & 0,9880 & 0,9880 & 17,76 & 0,365 \\
bilstm & glove    & 0,9877 & 0,9879 & 0,9877 & 0,9877 & 39,21 & 0,633 \\
lstm   & glove    & 0,9854 & 0,9855 & 0,9854 & 0,9854 & 55,69 & 0,680 \\
lstm   & word2vec & 0,9844 & 0,9846 & 0,9844 & 0,9844 & 64,12 & 0,415 \\
\bottomrule
\end{tabular}
\end{table}

Wszystkie architektury głębokie na ISOT przekraczają \(F_1 = 0{,}98\),
a~najlepsza konfiguracja (GRU + Word2Vec) osiąga \(F_1 = 0{,}9946\).
Średnia \(F_1 = 0{,}9896\) jest najwyższa spośród wszystkich rodzin modeli
z~wyjątkiem transformerów (por.\ Tabela~\ref{tab:bert_isot}). Word2Vec
systematycznie przewyższa GloVe -- trzy pierwsze miejsca rankingu należą
do kombinacji z~tym osadzeniem. Architektury rekurencyjne (GRU, BiLSTM)
nieznacznie wyprzedzają CNN, co można wiązać z~koniecznością modelowania
długich zależności sekwencyjnych w~pełnych artykułach prasowych.

\begin{table}[ht]
\centering
\caption{Wyniki sieci głębokich na zbiorze testowym WELFake.}
\label{tab:dl_welfake}
\begin{tabular}{llrrrrrr}
\toprule
Model & Embedding & Acc. & Prec. & Recall & $F_1$ & $t_{tr}$ [s] & $t_{inf}$ [s] \\
\midrule
gru    & word2vec & 0,9824 & 0,9824 & 0,9824 & 0,9824 & 144,14 & 0,999 \\
bilstm & word2vec & 0,9812 & 0,9812 & 0,9812 & 0,9812 & 166,63 & 1,267 \\
cnn    & word2vec & 0,9788 & 0,9788 & 0,9788 & 0,9788 & 129,05 & 0,825 \\
gru    & glove    & 0,9783 & 0,9783 & 0,9783 & 0,9783 & 48,71  & 1,008 \\
cnn    & glove    & 0,9769 & 0,9771 & 0,9769 & 0,9769 & 36,25  & 0,798 \\
bilstm & glove    & 0,9761 & 0,9762 & 0,9761 & 0,9761 & 71,27  & 1,257 \\
lstm   & glove    & 0,9669 & 0,9670 & 0,9669 & 0,9669 & 49,79  & 1,023 \\
lstm   & word2vec & 0,9604 & 0,9605 & 0,9604 & 0,9604 & 193,55 & 1,306 \\
\bottomrule
\end{tabular}
\end{table}

Na WELFake sieci głębokie utrzymują przewagę nad klasycznymi modelami:
średnia \(F_1 = 0{,}9751\) oznacza zysk ok.\ 9~pp.\ wobec tej rodziny.
Najsłabszą architekturą jest jednokierunkowy LSTM -- brak kontekstu
,,z~przyszłości'' ogranicza modelowanie narracyjnych struktur tekstu.
BiLSTM i~GRU osiągają zbliżone wyniki, lecz GRU trenuje się wyraźnie
szybciej.

Odmiennie kształtują się wyniki na zbiorze LIAR (Tabela~\ref{tab:dl_liar}).

\begin{table}[ht]
\centering
\caption{Wyniki sieci głębokich na zbiorze testowym LIAR.}
\label{tab:dl_liar}
\begin{tabular}{llrrrrr}
\toprule
Model & Embedding & Acc. & Prec. & Recall & $F_1$ & Uwagi \\
\midrule
bilstm & glove    & 0,6179 & 0,6152 & 0,6179 & 0,6157 & --- \\
cnn    & glove    & 0,6092 & 0,6040 & 0,6092 & 0,6006 & --- \\
cnn    & word2vec & 0,5638 & 0,5576 & 0,5638 & 0,5580 & --- \\
bilstm & word2vec & 0,5583 & 0,5795 & 0,5583 & 0,5562 & --- \\
gru    & word2vec & 0,5583 & 0,3117 & 0,5583 & 0,4001 & klasa większościowa \\
gru    & glove    & 0,4417 & 0,1951 & 0,4417 & 0,2706 & klasa większościowa \\
lstm   & word2vec & 0,4417 & 0,1951 & 0,4417 & 0,2706 & klasa większościowa \\
lstm   & glove    & 0,4417 & 0,1951 & 0,4417 & 0,2706 & klasa większościowa \\
\bottomrule
\end{tabular}
\end{table}

Wyniki LSTM i~GRU na LIAR jednoznacznie wskazują na zjawisko \textit{collapse}:
macierze pomyłek $\bigl(\begin{smallmatrix}564 & 0\\ 713 & 0\end{smallmatrix}\bigr)$
i~$\bigl(\begin{smallmatrix}0 & 564\\ 0 & 713\end{smallmatrix}\bigr)$
oznaczają, że sieci przewidywały wyłącznie jedną klasę. Wynika to
z~połączenia bardzo krótkich sekwencji (mediana ok.\ 18 słów) i~małego
zbioru treningowego (10~269 wypowiedzi) -- warunki niewystarczające dla
głębokich modeli rekurencyjnych. Spośród czterech architektur jedynie
BiLSTM z~GloVe i~CNN potrafią przekroczyć baseline większościowy; najlepszy
z~nich, BiLSTM + GloVe, osiąga \(F_1 = 0{,}6157\) -- wynik porównywalny
z~najlepszymi klasycznymi modelami na tym zbiorze.

\subsection{Analiza czasów uczenia i~wnioskowania}
\label{sec:wyniki_glebokie_czas}

\begin{table}[ht]
\centering
\caption{Średnie czasy uczenia i~wnioskowania sieci głębokich
w~zależności od zbioru danych (w~sekundach).}
\label{tab:dl_czasy}
\begin{tabular}{lrrrr}
\toprule
Zbiór & $\overline{t_{tr}}$ & $\max{t_{tr}}$ & $\overline{t_{inf}}$ & $\max{t_{inf}}$ \\
\midrule
ISOT    & 56,66  & 83,04  & 0,509 & 0,680 \\
LIAR    & 9,43   & 29,45  & 0,079 & 0,220 \\
WELFake & 104,92 & 193,55 & 1,060 & 1,306 \\
\bottomrule
\end{tabular}
\end{table}

Czasy uczenia sieci głębokich są 3--5-krotnie dłuższe niż klasycznych
modeli (np.\ \(104{,}9\)~s vs.\ \(33{,}3\)~s dla WELFake), lecz o~rząd
wielkości krótsze od transformerów. Czas inferencji na pełnym zbiorze
testowym mieści się w~granicach sekundy, co czyni tę rodzinę akceptowalną
w~zastosowaniach czasu rzeczywistego.

\subsection{Podsumowanie}
\label{sec:wyniki_glebokie_podsum}

Sieci głębokie zajmują pośrednią pozycję między klasycznymi algorytmami
a~transformerami. Na ISOT i~WELFake przekraczają \(F_1 = 0{,}97\), a~GRU + Word2Vec jest niemal
równorzędny z~XGBoost + TF-IDF. Na LIAR ujawniają się fundamentalne
ograniczenia tej rodziny: modele rekurencyjne w~niekorzystnych warunkach
mogą całkowicie zapaść na klasę większościową. Zdolność generalizacyjna
sieci głębokich poza dystrybucję treningową omówiona jest w~sekcji~%
\ref{sec:wyniki_www}.

\section{Wyniki modeli opartych na transformerach}
\label{sec:wyniki_transformery}

W~rozdziale tym omówiono wyniki czterech modeli z~rodziny Transformer:
BERT (\textit{bert-base-uncased}), DistilBERT (\textit{distilbert-base-uncased}),
RoBERTa (\textit{roberta-base}) oraz FakeBERT (wariant bazujący na
\textit{bert-base-uncased}). Każdy model dostrojono (\textit{fine-tuning})
osobno na każdym ze zbiorów z~użyciem optymalizatora AdamW i~schedulera
z~liniowym rozgrzewaniem. Learning rate, rozmiar batcha i~liczbę
epok dobrano za pomocą Optuna z~\textit{MedianPruner}.

\subsection{Wyniki na zbiorach testowych}
\label{sec:wyniki_bert_test}

\begin{table}[ht]
\centering
\caption{Wyniki modeli transformerowych na zbiorze testowym ISOT.}
\label{tab:bert_isot}
\begin{tabular}{llrrrrrr}
\toprule
Model & Tokenizer & Acc. & Prec. & Recall & $F_1$ & $t_{tr}$ [s] & $t_{inf}$ [s] \\
\midrule
fakebert   & bert-base-uncased       & 1,0000 & 1,0000 & 1,0000 & 1,0000 & 302,70 & 3,54 \\
roberta    & roberta-base            & 0,9995 & 0,9995 & 0,9995 & 0,9995 & 179,61 & 1,81 \\
bert       & bert-base-uncased       & 0,9990 & 0,9990 & 0,9990 & 0,9990 & 264,51 & 3,03 \\
distilbert & distilbert-base-uncased & 0,9982 & 0,9982 & 0,9982 & 0,9982 & 106,44 & 1,42 \\
\bottomrule
\end{tabular}
\end{table}

Na ISOT wszystkie transformery osiągnęły wyniki bliskie ideałowi.
FakeBERT uzyskał \(F_1 = 1{,}0000\) -- spośród 3911 próbek testowych
nie popełnił ani jednego błędu (macierz pomyłek
$\bigl(\begin{smallmatrix}1791 & 0\\ 0 & 2120\end{smallmatrix}\bigr)$).
Pozostałe modele przekroczyły barierę \(F_1 = 0{,}998\). Warto jednak
zwrócić uwagę, że na tym zbiorze przewaga transformerów nad klasycznymi
modelami wynosi jedynie 0{,}2--0{,}3~pp.\ \(F_1\), podczas gdy ich koszt
obliczeniowy jest ok.\ 10-krotnie wyższy -- co rodzi pytanie o~opłacalność
ich stosowania.

\begin{table}[ht]
\centering
\caption{Wyniki modeli transformerowych na zbiorze testowym WELFake.}
\label{tab:bert_welfake}
\begin{tabular}{llrrrrrr}
\toprule
Model & Tokenizer & Acc. & Prec. & Recall & $F_1$ & $t_{tr}$ [s] & $t_{inf}$ [s] \\
\midrule
roberta    & roberta-base            & 0,9978 & 0,9978 & 0,9978 & 0,9978 & 702,81 & 5,79  \\
bert       & bert-base-uncased       & 0,9951 & 0,9951 & 0,9951 & 0,9951 & 563,95 & 5,54  \\
fakebert   & bert-base-uncased       & 0,9950 & 0,9950 & 0,9950 & 0,9950 & 856,57 & 10,41 \\
distilbert & distilbert-base-uncased & 0,9943 & 0,9944 & 0,9943 & 0,9943 & 286,44 & 3,64  \\
\bottomrule
\end{tabular}
\end{table}

Bardziej wymagający WELFake uwidacznia wymierną przewagę transformerów.
RoBERTa osiągnęła tu \(F_1 = 0{,}9978\), co oznacza zysk ok.\ 2{,}6~pp.\
wobec najlepszego klasycznego modelu (XGBoost + BoW) i~1{,}5~pp.\ wobec
GRU + Word2Vec. Wszystkie cztery transformery mieszczą się w~wąskim
paśmie \(0{,}9943\)--\(0{,}9978\), potwierdzając wewnętrzną spójność
podejścia opartego na mechanizmie uwagi.

\begin{table}[ht]
\centering
\caption{Wyniki modeli transformerowych na zbiorze testowym LIAR.}
\label{tab:bert_liar}
\begin{tabular}{llrrrrrr}
\toprule
Model & Tokenizer & Acc. & Prec. & Recall & $F_1$ & $t_{tr}$ [s] & $t_{inf}$ [s] \\
\midrule
roberta    & roberta-base            & 0,6578 & 0,6560 & 0,6578 & 0,6564 & 50,48  & 0,400 \\
bert       & bert-base-uncased       & 0,6531 & 0,6537 & 0,6531 & 0,6534 & 31,98  & 0,685 \\
distilbert & distilbert-base-uncased & 0,6421 & 0,6434 & 0,6421 & 0,6426 & 27,95  & 0,263 \\
fakebert   & bert-base-uncased       & 0,6406 & 0,6380 & 0,6406 & 0,6383 & 285,38 & 1,349 \\
\bottomrule
\end{tabular}
\end{table}

Na LIAR transformery uzyskały najwyższe wyniki w~całym badaniu
(RoBERTa: \(F_1 = 0{,}6564\)). Choć przewaga nad najlepszym
klasycznym modelem (MNB + BoW: \(F_1 = 0{,}6326\)) wynosi jedynie
ok.\ 2{,}4~pp., jest ona istotnie większa niż w~przypadku sieci
głębokich, z~których część w~ogóle nie przekroczyła baseline'u
większościowego. FakeBERT, mimo najdłuższego czasu treningu spośród
transformerów, okazał się najsłabszym modelem tej rodziny --
pretrening na dłuższych artykułach nie przekłada się na lepszą
klasyfikację krótkich wypowiedzi politycznych, a~wręcz utrudnia
adaptację do tej domeny.

\subsection{Analiza czasów uczenia i~wnioskowania}
\label{sec:wyniki_bert_czas}

\begin{table}[ht]
\centering
\caption{Średnie czasy uczenia i~wnioskowania transformerów
w~zależności od zbioru danych (w~sekundach).}
\label{tab:bert_czasy}
\begin{tabular}{lrrrr}
\toprule
Zbiór & $\overline{t_{tr}}$ & $\max{t_{tr}}$ & $\overline{t_{inf}}$ & $\max{t_{inf}}$ \\
\midrule
ISOT    & 213,32 & 302,70 & 2,453 & 3,544  \\
LIAR    &  98,95 & 285,38 & 0,674 & 1,349  \\
WELFake & 602,44 & 856,57 & 6,346 & 10,411 \\
\bottomrule
\end{tabular}
\end{table}

Transformery są zdecydowanie najbardziej zasobochłonną rodziną modeli.
Średni czas uczenia na WELFake wynosi \(602{,}4\)~s -- ok.\ 18-krotnie
więcej niż dla klasycznych algorytmów i~6-krotnie więcej niż dla sieci
głębokich. Wnioskowanie zajmuje 1--10~ms na próbkę, co przy batchu 32
przekłada się na przepustowość ok.\ 200~próbek/s na GPU klasy T4.
Najszybszym transformerem jest DistilBERT, oferujący ok.\ 40\% niższy
czas inferencji niż BERT przy zbliżonej skuteczności.

\subsection{Podsumowanie}
\label{sec:wyniki_bert_podsum}

Transformery osiągają najwyższe wyniki we wszystkich trzech zbiorach.
Na ISOT doskonałą skuteczność uzyskuje FakeBERT (\(F_1 = 1{,}0000\)),
na WELFake dominuje RoBERTa (\(F_1 = 0{,}9978\)), a~na najtrudniejszym
LIAR ponownie RoBERTa (\(F_1 = 0{,}6564\)) -- choć przewaga nad
klasycznymi modelami jest tu umiarkowana. Główną słabością tej rodziny
pozostaje wysoki koszt obliczeniowy: kilkukrotnie wyższy niż pozostałych
rodzin zarówno na etapie uczenia, jak i~wnioskowania. Skuteczność
transformerów poza dystrybucją treningową omówiono w~sekcji~%
\ref{sec:wyniki_www} -- wykazano tam, że ich przewaga nad sieciami
głębokimi zanika lub odwraca się.

\section{Wyniki na danych pozyskanych ze stron WWW}
\label{sec:wyniki_www}

Rozdział ten przedstawia wyniki ewaluacji wszystkich trzech rodzin
modeli na zewnętrznym zbiorze artykułów pobranych metodą web-scrapingu
z~publicznie dostępnych portali informacyjnych i~serwisów satyrycznych.
Celem jest ocena zdolności do generalizacji poza dystrybucję treningową
-- kluczowego kryterium przydatności w~zastosowaniach produkcyjnych,
takich jak rozszerzenie przeglądarki.

\subsection{Charakterystyka zbioru ewaluacyjnego}
\label{sec:wyniki_www_dane}

Zewnętrzny zbiór ewaluacyjny pobrano automatycznie z~publicznie dostępnych
kanałów RSS i~stron internetowych. Podzielono go na dwie klasy: artykuły
z~uznanych portali informacyjnych (\emph{prawdziwe}, $y = 1$) oraz
publikacje serwisów satyrycznych i~parodystycznych (\emph{fałszywe},
$y = 0$). Wybór treści satyrycznych jako klasy negatywnej wynika
z~ich podobieństwa stylistycznego i~strukturalnego do rzetelnych
artykułów prasowych, co sprawia, że zadanie klasyfikacji jest
trudniejsze i~bliższe realnemu scenariuszowi niż użycie tekstów
generowanych syntetycznie.

Listę kandydujących źródeł RSS zaczerpnięto z~dwóch katalogów
prowadzonych przez serwis Feedspot: \textit{World News RSS Feeds}~\cite{feedspot_worldnews}
i~\textit{Fake News Websites RSS Feeds}~\cite{feedspot_fakenews}.
Po weryfikacji dostępności kanałów, języka i~częstotliwości aktualizacji
do dalszej pracy zakwalifikowano 50 serwisów informacyjnych i~12 serwisów
satyrycznych. Pełna lista źródeł i~selektory CSS do ekstrakcji treści
zdefiniowane są w~module \texttt{research/web/sources\_config.py}.

\subsubsection*{Źródła klasy pozytywnej (prawdziwe informacje)}

Klasę pozytywną stanowiło 50 kanałów RSS pochodzących z~globalnych
i~regionalnych serwisów informacyjnych. Tabela~\ref{tab:zrodla_prawda}
zawiera zestawienie źródeł pogrupowanych według regionu
geograficznego.

\begin{table}[ht]
\centering
\caption{Źródła klasy pozytywnej (prawdziwe informacje)
wykorzystane w~ewaluacji modeli.}
\label{tab:zrodla_prawda}
\begin{tabular}{p{3.2cm}p{10.5cm}}
\toprule
Region & Serwisy informacyjne \\
\midrule
USA i~Kanada & The New York Times, NPR, The Washington Post,
The Wall Street Journal, ABC News, CBS News, NBC News, Fox News,
NY Post, Time, CNBC, Los Angeles Times, Scripps News,
The Christian Science Monitor, The Globe and Mail, Global News \\
\midrule
Wielka Brytania & BBC World, BBC Tech, The Guardian (World, US),
Sky News, Daily Mail, The Independent, Express, Financial Times,
Daily Mirror, Daily Express \\
\midrule
Europa kontynentalna & Der Spiegel International, RFI \\
\midrule
Azja i~Bliski Wschód & Al Jazeera, The Print, NDTV, Tribune India,
The Hindu Businessline, The Week, CNBC TV India, Patrika News,
IFP (Iran), Morocco World News \\
\midrule
Australia i~Oceania & SBS News, The Sydney Morning Herald, 9News \\
\midrule
Inne / globalne & BuzzFeed World, Breaking News, Defence Blog,
Fair Observer, USNN \\
\bottomrule
\end{tabular}
\end{table}

\subsubsection*{Źródła klasy negatywnej (informacje fałszywe /
satyryczne)}

Klasę negatywną stanowiło 12 anglojęzycznych serwisów publikujących
treści fałszywe lub satyryczne, w~tym jedne z~najbardziej znanych
portali tego typu (\textit{The Onion}, \textit{The Babylon Bee}).
Lista została przedstawiona w~Tabeli~\ref{tab:zrodla_falsz}.

\begin{table}[ht]
\centering
\caption{Źródła klasy negatywnej (treści fałszywe i~satyryczne)
wykorzystane w~ewaluacji modeli.}
\label{tab:zrodla_falsz}
\begin{tabular}{p{3.2cm}p{10.5cm}}
\toprule
Region & Serwisy satyryczne / fałszywe \\
\midrule
USA & The Onion, The Babylon Bee, Duffel Blog (satyra wojskowa),
Gomer Blog (satyra medyczna), Empire News, Glossy News,
Weekly World News, Fake Ballmer \\
\midrule
Wielka Brytania & The Daily Mash, NewsBiscuit, The Poke, The Spoof \\
\midrule
Kanada & The Beaverton \\
\midrule
Australia & The Chaser \\
\bottomrule
\end{tabular}
\end{table}

\subsubsection*{Uzasadnienie wyboru serwisów satyrycznych jako klasy negatywnej}

Satyra nie jest tożsama z~dezinformacją -- ma świadomy charakter komediowy,
a~jej autorzy nie zamierzają wprowadzać czytelnika w~błąd. Przyjęcie
treści satyrycznych jako klasy fałszywej jest świadomą decyzją projektową,
uzasadnioną następującymi argumentami:

\begin{enumerate}
  \item \textbf{Dostępność danych w~skali tysięcy artykułów.}
        Serwisy satyryczne regularnie publikują nowe treści przez kanały
        RSS i~jawnie deklarują fikcyjny charakter publikacji. Umożliwia
        to automatyczne, etyczne pozyskanie etykietowanych przykładów
        klasy negatywnej, czego nie sposób osiągnąć dla faktycznej
        dezinformacji politycznej bez kosztownej, ręcznej anotacji przez
        zespół ekspertów.
  \item \textbf{Wysokie podobieństwo stylistyczne do rzetelnych artykułów.}
        W~odróżnieniu od tekstów syntetycznych, publikacje satyryczne
        celowo naśladują konwencje dziennikarskie: zachowują strukturę
        nagłówek-lead-rozwinięcie, stosują cytowania i~nazewnictwo osób.
        Czyni to zadanie klasyfikacji trudniejszym i~bliższym scenariuszowi
        rzeczywistemu, w~którym dezinformacja też imituje styl rzetelnych mediów.
  \item \textbf{Spójność z~dystrybucją treningową.} Wszystkie trzy zbiory
        treningowe (ISOT, LIAR, WELFake) zawierają heterogeniczne treści
        ,,fałszywe'' -- od dezinformacji politycznej, przez materiały
        spiskowe, po teksty satyryczne (szczególnie WELFake, łączący cztery
        podzbiory). Ewaluacja na serwisach satyrycznych jest spójna
        z~tym charakterem i~nie wprowadza sygnałów nieznanych modelowi.
  \item \textbf{Praktyczne znaczenie detekcji satyry.} Treści satyryczne
        są w~Internecie nierzadko odbierane jako prawdziwe -- szczególnie
        gdy udostępnia się sam nagłówek bez kontekstu źródła. Detekcja
        satyry to realne, dobrze udokumentowane zadanie, stanowiące
        prostszy wariant detekcji dezinformacji.
\end{enumerate}

Należy jednocześnie zaznaczyć ograniczenie tej metodologii: modele
uczą się cech charakterystycznych satyry (absurd, hiperbola,
przewidywalność puenty), nie zaś intencji manipulacyjnej właściwej
rzeczywistej dezinformacji. Wnioski należy interpretować w~kontekście
tej operacjonalizacji; bezpośrednie wdrożenie w~detekcji dezinformacji
politycznej wymagałoby walidacji na zbiorze faktycznie manipulacyjnych treści.

\subsubsection*{Metodologia pobierania danych}

Artykuły pobierano automatycznie skryptem \texttt{collect\_web.py},
który korzysta z~kanałów RSS i~selektorów CSS dopasowanych do struktury
HTML każdego serwisu. Dla stron nieujętych w~konfiguracji szczegółowej
(\texttt{SITE\_SELECTORS}) stosowano generyczne selektory standardowych
znaczników artykułów (\texttt{article p}, \texttt{div.entry-content p}).
W~celu ograniczenia ryzyka blokad nagłówki \texttt{User-Agent} losowano
spośród trzech przeglądarek desktopowych. Wszystkie kanały RSS
zostały zweryfikowane pod kątem aktualności, języka i~spójności
tematycznej z~zadaniem klasyfikacji.

Dla każdego artykułu wyodrębniono trzy warianty tekstowe:
\begin{itemize}
  \item \textbf{title} -- tytuł artykułu (typowo do 30 słów),
  \item \textbf{short text} -- pierwsze 1--3 akapity lub podsumowanie (30--100 słów),
  \item \textbf{text} -- pełna treść (powyżej 100 słów).
\end{itemize}

Zebrano łącznie 1599 próbek dla wariantu \emph{text},
1590 dla \emph{title} i~1593 dla \emph{short text}. Zbiór jest
istotnie niezbalansowany: stosunek klasy pozytywnej do negatywnej
wynosi ok.\ \(9{:}1\) (1436 artykułów prawdziwych vs.\ 163 satyrycznych
dla wariantu \emph{text}). Wynika to z~ograniczonej aktywności
kanałów RSS serwisów satyrycznych. Z~tego powodu w~prezentowanych
wynikach koncentrujemy się na precyzji i~czułości w~rozbiciu na klasy,
a~nie wyłącznie na ogólnej dokładności, która mogłaby być zawyżona
przez baseline większościowy.

\subsection{Wyniki klasycznych modeli uczenia maszynowego}
\label{sec:wyniki_www_klasyczne}

Tabela~\ref{tab:klas_web} zawiera najlepsze wyniki klasycznych
modeli w~każdym z~trzech wariantów długości tekstu.

\begin{table}[ht]
\centering
\caption{Najlepsze klasyczne modele w~ewaluacji na danych WWW
(według miary $F_1$).}
\label{tab:klas_web}
\begin{tabular}{lllrrrr}
\toprule
Wariant & Zbiór & Model + Emb. & Acc. & Prec. & Recall & $F_1$ \\
\midrule
text       & LIAR    & dt + tfidf      & 0,8074 & 0,8210 & 0,8074 & 0,8140 \\
text       & ISOT    & nb + word2vec   & 0,6929 & 0,8969 & 0,6929 & 0,7533 \\
text       & WELFake & dt + bow        & 0,7136 & 0,7867 & 0,7136 & 0,7483 \\
title      & ISOT    & rf + bow        & 0,8799 & 0,8350 & 0,8799 & 0,8524 \\
title      & WELFake & rf + tfidf      & 0,8981 & 0,8066 & 0,8981 & 0,8499 \\
title      & LIAR    & svm + tfidf     & 0,6201 & 0,7800 & 0,6201 & 0,6934 \\
short text & ISOT    & rf + bow        & 0,8970 & 0,8684 & 0,8970 & 0,8747 \\
short text & WELFake & dt + bow        & 0,8839 & 0,8119 & 0,8839 & 0,8451 \\
short text & LIAR    & svm + tfidf     & 0,6554 & 0,7990 & 0,6554 & 0,7210 \\
\bottomrule
\end{tabular}
\end{table}

Klasyczne modele zaskakująco dobrze generalizują poza zbiór
treningowy, szczególnie dla nagłówków i~leadów. Najlepsze konfiguracje
(RF + BoW, DT + BoW) na tych wariantach uzyskują
\(F_1 \in [0{,}84;\,0{,}87]\) -- nieznacznie poniżej wyników in-domain.
Na pełnych tekstach średnia spada do ok.\ \(0{,}58\), co świadczy
o~silnej zależności klasycznych reprezentacji od częstości konkretnych
słów -- co ogranicza generalizację między stylistycznie odmiennymi
korpusami.

\subsection{Wyniki sieci głębokich}
\label{sec:wyniki_www_glebokie}

Wyniki sieci głębokich na danych WWW zostały zestawione
w~Tabeli~\ref{tab:dl_web}.

\begin{table}[ht]
\centering
\caption{Najlepsze sieci głębokie w~ewaluacji na danych WWW
(według miary $F_1$).}
\label{tab:dl_web}
\begin{tabular}{lllrrrr}
\toprule
Wariant & Zbiór & Model + Emb. & Acc. & Prec. & Recall & $F_1$ \\
\midrule
text       & LIAR    & bilstm + glove    & 0,8843 & 0,8388 & 0,8843 & 0,8550 \\
text       & WELFake & bilstm + word2vec & 0,5935 & 0,8073 & 0,5935 & 0,6732 \\
text       & ISOT    & cnn + word2vec    & 0,5879 & 0,8826 & 0,5879 & 0,6671 \\
title      & LIAR    & gru + word2vec    & 0,8981 & 0,8066 & 0,8981 & 0,8499 \\
title      & WELFake & lstm + word2vec   & 0,8874 & 0,8162 & 0,8874 & 0,8467 \\
title      & ISOT    & cnn + word2vec    & 0,5434 & 0,8800 & 0,5434 & 0,6274 \\
short text & LIAR    & gru + word2vec    & 0,8977 & 0,8090 & 0,8977 & 0,8511 \\
short text & WELFake & lstm + word2vec   & 0,8286 & 0,8158 & 0,8286 & 0,8221 \\
short text & ISOT    & cnn + word2vec    & 0,6321 & 0,8800 & 0,6321 & 0,7050 \\
\bottomrule
\end{tabular}
\end{table}

Sieci głębokie na danych WWW wykazują interesującą prawidłowość:
najlepiej generalizują modele trenowane na LIAR i~WELFake, mimo
że in-domain plasowały się najniżej. Najwyższy wynik w~całej
ewaluacji web osiąga BiLSTM + GloVe trenowany na LIAR
(\(F_1 = 0{,}8550\)) dla pełnych tekstów -- wyrównując najlepsze
klasyczne modele i~wyprzedzając transformery. Modele trenowane
na ISOT tracą natomiast jakość na danych zewnętrznych
(\(F_1\) zaledwie \(0{,}63\)--\(0{,}71\)), co potwierdza obecność
silnych korelacji stylistycznych w~tym korpusie (prefiksy ,,Reuters'',
stałe sygnatury) -- korelacji, które wyciekają do zbioru testowego
in-domain.

\subsection{Wyniki modeli opartych na transformerach}
\label{sec:wyniki_www_transformery}

Tabela~\ref{tab:bert_web} prezentuje najlepsze wyniki transformerów
na danych WWW.

\begin{table}[ht]
\centering
\caption{Najlepsze modele transformerowe w~ewaluacji na danych WWW
(według miary $F_1$).}
\label{tab:bert_web}
\begin{tabular}{lllrrrr}
\toprule
Wariant & Zbiór & Model & Acc. & Prec. & Recall & $F_1$ \\
\midrule
text       & ISOT    & roberta    & 0,7311 & 0,8288 & 0,7311 & 0,7724 \\
text       & WELFake & fakebert   & 0,6629 & 0,7952 & 0,6629 & 0,7205 \\
text       & LIAR    & fakebert   & 0,5641 & 0,8247 & 0,5641 & 0,6498 \\
title      & ISOT    & roberta    & 0,7465 & 0,8630 & 0,7465 & 0,7899 \\
title      & WELFake & fakebert   & 0,7547 & 0,7960 & 0,7547 & 0,7747 \\
title      & LIAR    & fakebert   & 0,2711 & 0,7905 & 0,2711 & 0,3273 \\
short text & ISOT    & roberta    & 0,8418 & 0,8741 & 0,8418 & 0,8556 \\
short text & ISOT    & fakebert   & 0,7571 & 0,8909 & 0,7571 & 0,8017 \\
short text & LIAR    & fakebert   & 0,5191 & 0,8622 & 0,5191 & 0,6078 \\
\bottomrule
\end{tabular}
\end{table}

Na danych WWW transformery prezentują niejednoznaczny obraz. Dla pełnych
tekstów osiągają najwyższą średnią \(F_1\) spośród wszystkich rodzin
(\(\overline{F_1} = 0{,}651\) vs.\ \(0{,}580\) dla klasycznych i~\(0{,}576\)
dla sieci głębokich). Jednak dla nagłówków i~krótkich fragmentów ulegają
sieciom głębokim trenowanym na LIAR. Najlepszą pojedynczą konfiguracją
jest RoBERTa trenowana na ISOT, która dla wariantu \emph{short text}
uzyskuje \(F_1 = 0{,}8556\) -- najwyższą wartość w~tej kategorii.

\subsection{Podsumowanie ewaluacji na danych WWW}
\label{sec:wyniki_www_podsum}

Ewaluacja na danych zewnętrznych ujawniła kilka kluczowych
prawidłowości niedostrzegalnych przy miarach in-domain:

\begin{itemize}
  \item \textbf{Wyraźny spadek skuteczności wszystkich rodzin.}
        W~porównaniu ze zbiorami testowymi miara $F_1$ spada
        średnio o~20--30~pp. Wszystkie trzy korpusy (ISOT, LIAR,
        WELFake) zawierają źródłospecyficzne sygnały dystrybucyjne,
        które nie generalizują poza własną pulę redakcyjną.
  \item \textbf{Najlepsze modele in-domain nie są najlepszymi
        modelami na danych WWW.} BiLSTM + GloVe trenowany na LIAR
        (\(F_1^{\text{in-domain}} = 0{,}6157\)) okazał się
        najlepszym klasyfikatorem pełnych tekstów WWW
        (\(F_1 = 0{,}8550\)). Stanowi to ważne odkrycie
        metodologiczne -- dobór modelu wyłącznie na podstawie
        miary in-domain prowadziłby do błędnych wniosków produkcyjnych.
  \item \textbf{Transformery tracą przewagę poza in-domain.}
        Dla nagłówków i~krótkich fragmentów transformery są często
        gorsze od klasycznych modeli i~sieci głębokich, przy
        znacznie wyższym koszcie obliczeniowym.
  \item \textbf{Optymalna rodzina modeli zależy od długości tekstu.}
        RF + BoW dominuje dla krótkich tekstów (\textit{title},
        \textit{short text}), BiLSTM + GloVe -- dla pełnych artykułów.
        Stanowi to uzasadnienie dla mechanizmu \textit{smart routingu}
        omówionego w~sekcji~\ref{sec:podsum_rekomendacja}.
\end{itemize}

Powyższe obserwacje stanowią punkt wyjścia do analizy kompromisu
skuteczność--koszt inferencji w~kolejnym rozdziale.

\section{Kompromis pomiędzy skutecznością predykcji a~czasem inferencji}
\label{sec:wyniki_kompromis}

Poprzednie sekcje (\ref{sec:wyniki_klasyczne}--\ref{sec:wyniki_www})
analizowały wyniki poszczególnych rodzin modeli niezależnie.
Niniejszy rozdział łączy obie osie -- skuteczność i~koszt inferencji
-- odpowiadając na pytanie, jaki model wybrać w~praktycznym zastosowaniu
detekcji fake newsów w~czasie rzeczywistym.

\subsection{Średni koszt inferencji według rodziny modeli}
\label{sec:kompromis_srednie}

Tabela~\ref{tab:kompromis_srednie} zestawia średnią miarę $F_1$
i~czas inferencji dla każdej rodziny, uśrednione po wszystkich
trzech zbiorach (ISOT, LIAR, WELFake).

\begin{table}[ht]
\centering
\caption{Średnia jakość i~koszt inferencji w~podziale na rodziny
modeli (wszystkie wyniki na zbiorach testowych).}
\label{tab:kompromis_srednie}
\begin{tabular}{lrrrr}
\toprule
Rodzina & $\overline{F_1}$ & $\overline{t_{inf}}$ [s] & $\overline{t_{tr}}$ [s] & lat./próbkę [ms] \\
\midrule
Klasyczne    & 0,8084 & 1,101 & 18,17  & 0,011\,(med.) \\
Sieci głębokie & 0,8025 & 0,549 & 57,00  & 0,125\,(med.) \\
Transformery & 0,8808 & 3,158 & 304,90 & 0,673\,(med.) \\
\bottomrule
\end{tabular}
\end{table}

Kolumna \emph{lat./próbkę} zawiera medianę stosunku $t_{inf}/N$
dla wszystkich eksperymentów (N = 3911 dla ISOT, 6368 dla WELFake,
1277 dla LIAR), zapewniając miarę niezależną od rozmiaru korpusu.

Z~zestawienia wynika, że transformery oferują przeciętnie
\(7{,}2\)~pp.\ wyższe $F_1$ od modeli klasycznych, lecz kosztem
ok.\ \emph{60-krotnie} dłuższego czasu inferencji na próbkę. Sieci
głębokie zajmują pozycję pośrednią, ale ich średnie \(F_1 = 0{,}8025\)
jest o~\(0{,}6\)~pp.\ niższe od klasycznych -- co wyklucza je jako
oczywisty ,,kompromis'' bez dodatkowej kontekstowej analizy.

\subsection{Najlepsze konfiguracje według rodziny modeli}
\label{sec:kompromis_best}

Tabela~\ref{tab:kompromis_best} zestawia najlepsze konfiguracje
(wg $F_1$) z~każdej rodziny dla każdego ze zbiorów, wraz
z~czasami uczenia i~inferencji, umożliwiając podejmowanie
konkretnych decyzji inżynierskich.

\begin{table}[ht]
\centering
\caption{Najlepsze konfiguracje z~każdej rodziny modeli wraz
z~kosztami obliczeniowymi. Symbol \,ms/próbka odnosi się do
oszacowanej średniej latencji pojedynczego zapytania.}
\label{tab:kompromis_best}
\small
\setlength{\tabcolsep}{4pt}
\begin{tabular}{@{}lllrrrr@{}}
\toprule
Zbiór & Rodzina & Model + emb. & $F_1$ & $t_{tr}$ [s] & $t_{inf}$ [s] & ms/próbkę \\
\midrule
\multirow{3}{*}{ISOT}
 & Klasyczne     & xgb + tfidf            & 0,9977 & 150,05 & 0,072  & 0,018 \\
 & Sieci głębokie & gru + word2vec        & 0,9946 & 67,06  & 0,491  & 0,125 \\
 & Transformery  & fakebert + bert-base   & 1,0000 & 302,70 & 3,544  & 0,906 \\
\midrule
\multirow{3}{*}{LIAR}
 & Klasyczne     & mnb + bow              & 0,6326 & 0,01   & 0,008  & 0,006 \\
 & Sieci głębokie & bilstm + glove        & 0,6157 & 8,48   & 0,085  & 0,066 \\
 & Transformery  & roberta + roberta-base & 0,6564 & 50,48  & 0,400  & 0,313 \\
\midrule
\multirow{3}{*}{WELFake}
 & Klasyczne     & xgb + bow              & 0,9713 & 56,17  & 0,173  & 0,027 \\
 & Sieci głębokie & gru + word2vec        & 0,9824 & 144,14 & 0,999  & 0,157 \\
 & Transformery  & roberta + roberta-base & 0,9978 & 702,81 & 5,793  & 0,910 \\
\bottomrule
\end{tabular}
\end{table}

We wszystkich trzech zbiorach transformer prowadzi pod względem $F_1$,
jednak różnica względem najlepszego modelu klasycznego jest minimalna:
0{,}23~pp.\ na ISOT, 2{,}38~pp.\ na LIAR i~2{,}65~pp.\ na WELFake.
Te skromne przyrosty są kupowane za cenę \(50\)--\(120\)-krotnie
dłuższego czasu inferencji. W~większości zastosowań praktycznych
zysk staje się marginalny ekonomicznie -- koszt obliczeniowy rośnie
szybciej niż przyrost jakości.

\subsection{Granica Pareto}
\label{sec:kompromis_pareto}

W~celu rygorystycznej analizy kompromisu skuteczność--koszt
przeprowadzono analizę Pareto. Konfiguracja $A = (t_{inf}^A,\,F_1^A)$
\emph{dominuje} konfigurację $B = (t_{inf}^B,\,F_1^B)$, gdy
$t_{inf}^A \le t_{inf}^B$ i~$F_1^A \ge F_1^B$, przy czym co najmniej
jedna nierówność jest ostra. Granica Pareto to zbiór konfiguracji
niedominowanych przez żadną inną -- takich, że poprawa jednego kryterium
jest możliwa wyłącznie kosztem pogorszenia drugiego~\cite{pareto1906manuale,
miettinen1999nonlinear, deb2001multi}. W~kontekście uczenia maszynowego
analiza ta nabiera szczególnego znaczenia w~obliczu rosnących kosztów
obliczeniowych nowoczesnych architektur (\emph{Green AI})~\cite{schwartz2020green,
treviso2023efficient}.

Dla każdego zbioru wyznaczono granicę Pareto w~przestrzeni
$(\,t_{inf},\,F_1\,)$.

\begin{table}[ht]
\centering
\caption{Granica Pareto dla zbioru ISOT.}
\label{tab:pareto_isot}
\begin{tabular}{lllrrr}
\toprule
\# & Rodzina & Model + emb. & $F_1$ & $t_{inf}$ [s] & $t_{tr}$ [s] \\
\midrule
1 & Klasyczne   & dt + word2vec          & 0,9304 & 0,0105 & 5,95   \\
2 & Klasyczne   & dt + bow               & 0,9378 & 0,0153 & 10,76  \\
3 & Klasyczne   & mnb + bow              & 0,9553 & 0,0158 & 0,04   \\
4 & Klasyczne   & lr + word2vec          & 0,9813 & 0,0160 & 3,30   \\
5 & Klasyczne   & svm + word2vec         & 0,9836 & 0,0233 & 6,60   \\
6 & Klasyczne   & lr + tfidf             & 0,9921 & 0,0403 & 1,77   \\
7 & Klasyczne   & xgb + bow              & 0,9974 & 0,0634 & 12,14  \\
8 & Klasyczne   & xgb + tfidf            & 0,9977 & 0,0717 & 150,05 \\
9 & Transformer   & distilbert             & 0,9982 & 1,4226 & 106,44 \\
10& Transformer   & roberta                & 0,9995 & 1,8146 & 179,61 \\
11& Transformer   & fakebert + bert-base   & 1,0000 & 3,5443 & 302,70 \\
\bottomrule
\end{tabular}
\end{table}

\begin{table}[ht]
\centering
\caption{Granica Pareto dla zbioru WELFake.}
\label{tab:pareto_welfake}
\begin{tabular}{lllrrr}
\toprule
\# & Rodzina & Model + emb. & $F_1$ & $t_{inf}$ [s] & $t_{tr}$ [s] \\
\midrule
1 & Klasyczne   & dt + word2vec          & 0,8654 & 0,0107 & 8,84   \\
2 & Klasyczne   & lr + word2vec          & 0,9165 & 0,0160 & 5,55   \\
3 & Klasyczne   & dt + bow               & 0,9364 & 0,0206 & 20,21  \\
4 & Klasyczne   & xgb + tfidf            & 0,9711 & 0,0449 & 147,67 \\
5 & Klasyczne   & xgb + bow              & 0,9713 & 0,1728 & 56,17  \\
6 & Sieci głębokie   & cnn + glove            & 0,9769 & 0,7975 & 36,25  \\
7 & Sieci głębokie   & cnn + word2vec         & 0,9788 & 0,8251 & 129,05 \\
8 & Sieci głębokie   & gru + word2vec         & 0,9824 & 0,9993 & 144,14 \\
9 & Transformer   & distilbert             & 0,9943 & 3,6418 & 286,44 \\
10& Transformer   & bert                   & 0,9951 & 5,5361 & 563,95 \\
11& Transformer   & roberta                & 0,9978 & 5,7928 & 702,81 \\
\bottomrule
\end{tabular}
\end{table}

\begin{table}[ht]
\centering
\caption{Granica Pareto dla zbioru LIAR.}
\label{tab:pareto_liar}
\begin{tabular}{lllrrr}
\toprule
\# & Rodzina & Model + emb. & $F_1$ & $t_{inf}$ [s] & $t_{tr}$ [s] \\
\midrule
1 & Klasyczne   & dt + word2vec          & 0,5648 & 0,0080 & 0,65   \\
2 & Klasyczne   & mnb + bow              & 0,6326 & 0,0082 & 0,01   \\
3 & Transformer   & distilbert             & 0,6426 & 0,2633 & 27,95  \\
4 & Transformer   & roberta                & 0,6564 & 0,4000 & 50,48  \\
\bottomrule
\end{tabular}
\end{table}

Z~granicy Pareto wynikają trzy obserwacje:

\begin{itemize}
  \item Na ISOT przestrzeń ,,szybkie i~niemal tak samo dokładne''
        należy w~całości do klasycznych modeli (8 z~11 punktów
        granicy). Transformer jest uzasadniony dopiero przy
        wymaganiach $F_1 > 0{,}9977$.
  \item Na WELFake sieci głębokie zajmują trzy środkowe punkty granicy,
        oferując pośredni kompromis między szybkim XGBoostem
        a~kosztownymi transformerami. Zysk \(1{,}1\)~pp.\ $F_1$
        (z~\(0{,}9713\) do \(0{,}9824\)) kosztuje wzrost latencji
        o~rząd wielkości.
  \item Na LIAR sieci głębokie nie pojawiają się na granicy Pareto
        w~ogóle: BiLSTM + GloVe (\(F_1 = 0{,}6157\)) jest zarówno
        wolniejszy, jak i~mniej dokładny od klasycznego MNB + BoW,
        a~więc przez niego \emph{dominowany}. Potwierdza to wniosek
        z~rozdz.~\ref{sec:wyniki_glebokie_podsum}, że dla krótkich
        tekstów sieci rekurencyjne nie wnoszą wartości dodanej.
\end{itemize}

\subsection{Koszt uczenia}
\label{sec:kompromis_train}

Obok kosztu inferencji istotny jest jednorazowy nakład na etapie
treningu. Średnie czasy uczenia różnią się między rodzinami o~rząd
wielkości: 18~s dla klasycznych, 57~s dla sieci głębokich,
305~s dla transformerów (Tabela~\ref{tab:kompromis_srednie}).
W~skrajnych przypadkach (FakeBERT na WELFake) trening trwa ponad
14~minut, co oznacza realny koszt zasobów GPU.

Z~praktycznego punktu widzenia kluczowy jest stosunek przyrostu $F_1$
do przyrostu kosztu uczenia. Zastępując XGBoost + TF-IDF przez RoBERTa
na WELFake, uzyskujemy zysk \(2{,}65\)~pp.\ $F_1$ za \(8{,}1\times\)
dłuższy trening. W~scenariuszach wymagających częstego retreningu
(np.\ codzienne dotrenowanie na nowych artykułach) to kryterium
często decyduje o~wyborze architektury.

\subsection{Kompromis na danych pozyskanych ze stron WWW}
\label{sec:kompromis_web}

Dotychczasowa analiza dotyczyła zbiorów testowych in-domain.
W~zastosowaniach produkcyjnych decydująca jest jednak skuteczność
na świeżych danych. Tabela~\ref{tab:kompromis_web} zestawia najlepsze
konfiguracje na danych pobranych z~Internetu (por.\ sekcja~\ref{sec:wyniki_www}).

\begin{table}[ht]
\centering
\caption{Najlepsze konfiguracje w~ewaluacji na danych pozyskanych
ze stron WWW z~uwzględnieniem czasu inferencji.}
\label{tab:kompromis_web}
\begin{tabular}{llllrr}
\toprule
Wariant & Rodzina & Model (zbiór) & $F_1$ & $t_{inf}$ [s] \\
\midrule
text       & Sieci głębokie   & bilstm + glove (LIAR)       & 0,8550 & 0,188 \\
text       & Klasyczne   & dt + tfidf (LIAR)           & 0,8140 & 1,005 \\
text       & Transformer   & roberta (ISOT)              & 0,7724 & 1,592 \\
title      & Klasyczne   & rf + bow (ISOT)             & 0,8524 & 0,015 \\
title      & Sieci głębokie   & gru + word2vec (LIAR)       & 0,8499 & 0,036 \\
title      & Transformer   & roberta (ISOT)              & 0,7899 & 0,868 \\
short text & Klasyczne   & rf + bow (ISOT)             & 0,8747 & 0,025 \\
short text & Transformer   & roberta (ISOT)              & 0,8556 & 1,247 \\
short text & Sieci głębokie   & gru + word2vec (LIAR)       & 0,8511 & 0,038 \\
\bottomrule
\end{tabular}
\end{table}

Na danych WWW \emph{transformery tracą pozycję lidera}. Dla wariantów
\emph{title} i~\emph{short text} najlepsze wyniki osiąga las losowy
z~BoW, będący jednocześnie \(50\)--\(80\times\) szybszy. Dla pełnych
tekstów prowadzi BiLSTM + GloVe, łącząc najwyższe \(F_1\) z~latencją
poniżej 0{,}2~s. Płynie z~tego czytelna rekomendacja: w~zastosowaniach
produkcyjnych uzasadniony jest mechanizm rotacji modeli zależny od
długości tekstu wejściowego.

\subsection{Wnioski z~analizy kompromisu}
\label{sec:kompromis_wnioski}

\begin{itemize}
  \item Przy ograniczeniu latencji $\le 50$\,ms/zapytanie (typowym
        dla interfejsów webowych) Pareto-optymalne są wyłącznie modele
        klasyczne -- najczęściej XGBoost lub Random Forest z~reprezentacją rzadką.
  \item Transformer jest uzasadniony wyłącznie wtedy, gdy maksymalna
        jakość klasyfikacji jest priorytetem i~dostępna jest karta GPU
        zapewniająca inferencję w~granicach 5--10~ms na próbkę.
  \item Sieci głębokie pełnią pożyteczną rolę ,,środka'' jedynie dla
        dłuższych tekstów (ISOT, WELFake); dla krótkich (LIAR, nagłówki)
        są dominowane przez klasyczne modele.
  \item Największą stabilność na danych z~innej dystrybucji wykazują
        modele drzewiaste (RF, DT) z~BoW. Uzupełnione o~mechanizm
        routingu osiągają kompromis nieosiągalny dla żadnego
        pojedynczego modelu.
\end{itemize}

\section{Podsumowanie wyników eksperymentów}
\label{sec:wyniki_podsumowanie}

Niniejsza sekcja syntetyzuje wyniki z~sekcji
\ref{sec:wyniki_klasyczne}--\ref{sec:wyniki_kompromis}.
Przeprowadzono łącznie 105 eksperymentów ewaluacyjnych na zbiorach
testowych (23 konfiguracje klasyczne $\times$ 3 + 8 sieci głębokich
$\times$ 3 + 4 transformery $\times$ 3) oraz 312 ewaluacji na danych
pobranych metodą web-scrapingu w~trzech wariantach
(\emph{text}, \emph{title}, \emph{short text}).

\subsection{Najlepsze modele w~podziale na zbiory}
\label{sec:podsum_best}

Tabela~\ref{tab:podsum_best} zestawia najlepsze konfiguracje
uzyskane w~całym badaniu, niezależnie od rodziny modelu.

\begin{table}[ht]
\centering
\caption{Najlepsze konfiguracje w~całym badaniu (na zbiorach
testowych).}
\label{tab:podsum_best}
\begin{tabular}{lllrr}
\toprule
Zbiór & Najlepszy model & Rodzina & $F_1$ & Liczba próbek \\
\midrule
ISOT    & fakebert + bert-base   & Transformer & 1,0000 & 3911 \\
WELFake & roberta + roberta-base & Transformer & 0,9978 & 6368 \\
LIAR    & roberta + roberta-base & Transformer & 0,6564 & 1277 \\
\bottomrule
\end{tabular}
\end{table}

We wszystkich trzech zbiorach najwyższe wyniki osiągnęły transformery.
Warta podkreślenia jest istotna dysproporcja trudności między korpusami:
$F_1$ na LIAR jest o~ponad 34~pp.\ niższe niż na ISOT. Potwierdza to
obserwację z~literatury -- krótkie, merytorycznie złożone wypowiedzi
polityczne stanowią znacznie trudniejsze zadanie niż klasyfikacja
pełnych artykułów prasowych.

\subsection{Średnia jakość w~poszczególnych scenariuszach ewaluacji}
\label{sec:podsum_srednia}

Tabela~\ref{tab:podsum_srednia} prezentuje uśrednione miary $F_1$
w~podziale na rodzinę modelu i~rodzaj ewaluacji.

\begin{table}[ht]
\centering
\caption{Średnia miara $F_1$ w~różnych scenariuszach ewaluacji
(uśrednienie po wszystkich modelach i~zbiorach danej rodziny).}
\label{tab:podsum_srednia}
\begin{tabular}{lrrrr}
\toprule
Rodzina        & test     & web-text & web-title & web-short \\
\midrule
Klasyczne      & 0,8084 & 0,5802 & 0,5479 & 0,5926 \\
Sieci głębokie & 0,8025 & 0,5761 & 0,5720 & 0,5965 \\
Transformery   & 0,8808 & 0,6512 & 0,5022 & 0,5942 \\
\bottomrule
\end{tabular}
\end{table}

Zestawienie ujawnia wyraźną prawidłowość: \emph{przewaga transformerów
jest największa w~scenariuszu in-domain} (\(0{,}88\) vs.\ \(0{,}81\)
dla klasycznych) oraz \emph{dla pełnych tekstów z~Internetu}
(\(0{,}65\) vs.\ \(0{,}58\)). Natomiast
w~scenariuszach krótkich tekstów (\emph{title}, \emph{short text})
transformery tracą przewagę i~są przeciętnie słabsze od sieci głębokich
i~klasycznych modeli. Przyczyną może być typowa długość tekstów
treningowych (ISOT i~WELFake liczą setki słów na próbkę) -- modele
dostrojone na długich artykułach słabo ekstrapolują do krótkich nagłówków.

\subsection{Najważniejsze obserwacje}
\label{sec:podsum_obserwacje}

Z~całokształtu eksperymentów wynikają następujące wnioski.

\begin{enumerate}
  \item \textbf{ISOT zawiera silne sygnały dystrybucyjne}, które
        pozwalają nawet prostym klasyfikatorom uzyskać $F_1$
        powyżej \(0{,}99\). Sugeruje to obecność wycieku informacji
        (np.\ stałych prefiksów ,,Reuters'') i~uzasadnia zastosowanie
        dodatkowego mechanizmu czyszczenia danych,
        który jednak nie zmniejszył drastycznie miary $F_1$.
  \item \textbf{LIAR pozostaje fundamentalnym wyzwaniem.} Żaden
        z~32~testowanych modeli (klasyczne, głębokie, transformerowe)
        nie przekroczył $F_1 = 0{,}66$. Cztery konfiguracje sieci
        głębokich (LSTM, GRU) zawalały się do klasy większościowej,
        co świadczy o~zbyt małym rozmiarze zbioru w~stosunku do
        zdolności modelu.
  \item \textbf{WELFake jest najbardziej zrównoważonym zbiorem} pod
        względem informatywności i~trudności. To na nim najlepiej
        widać przewagę kolejnych rodzin: klasyczny XGBoost
        (\(0{,}9713\)) $\to$ GRU + Word2Vec (\(0{,}9824\)) $\to$
        RoBERTa (\(0{,}9978\)).
  \item \textbf{Reprezentacja BoW/TF-IDF dominuje} wśród klasycznych
        modeli na każdym zbiorze, podczas gdy \textbf{Word2Vec lepiej
        spisuje się dla sieci głębokich} -- ze względu na semantyczne
        sąsiedztwo wektorów.
  \item \textbf{Najlepsza generalizacja na dane WWW należy do
        BiLSTM~+~GloVe wytrenowanego na LIAR}, co stanowi
        nieoczekiwany rezultat: model najsłabszy in-domain okazał
        się najmocniejszy w~scenariuszu transferu domenowego.
  \item \textbf{Kompromis kosztu i~jakości} jest najbardziej
        opłacalny przy zastosowaniu klasycznych modeli z~rzadkimi
        reprezentacjami. Dla większości zastosowań produkcyjnych
        XGBoost lub RandomForest oferują skuteczność rzędu
        \(0{,}97\)--\(0{,}99\) przy latencji poniżej \(0{,}05\)~ms
        na próbkę.
\end{enumerate}

\subsection{Rekomendacja modelu dla wtyczki Chrome}
\label{sec:podsum_rekomendacja}

Docelowym scenariuszem zastosowania wytrenowanych modeli jest
wtyczka do przeglądarki Chrome analizująca zaznaczony przez
użytkownika fragment artykułu. W~praktyce długość takiego fragmentu
może być bardzo różna: od kilku słów nagłówka, przez krótki opis
(lead), aż po pełny tekst artykułu. Z~tego powodu projekt zakłada
mechanizm \emph{smart routingu} -- selekcji modelu w~zależności od
długości tekstu wejściowego.

W~niniejszej podsekcji przedstawiono niezależnie wyprowadzoną
rekomendację trzech najlepszych konfiguracji dla każdego z~trzech
zakresów długości tekstu, posługując się następującymi kryteriami:

\begin{itemize}
  \item $F_1^{\text{web}}$ -- jakość klasyfikacji na danych
        pozyskanych ze stron WWW (najwyższy priorytet, ponieważ
        odzwierciedla rzeczywiste warunki pracy wtyczki),
  \item $F_1^{\text{test}}$ -- jakość klasyfikacji na zbiorze
        testowym in-domain (informacja o~stabilności modelu),
  \item $t_{\text{lat}}$ -- mediana latencji pojedynczej predykcji
        wyrażona w~milisekundach (kluczowa dla interfejsu czasu
        rzeczywistego),
  \item $Q = 0{,}6 \cdot F_1^{\text{web}} + 0{,}4 \cdot F_1^{\text{test}}$
        -- syntetyczna miara łącząca obie jakości, z~wagą
        faworyzującą skuteczność na danych spoza dystrybucji
        treningowej.
\end{itemize}

\subsubsection*{Rekomendacja dla nagłówków (\textit{title}, $<30$ słów)}

\begin{table}[ht]
\centering
\caption{Top~3 modele dla klasyfikacji nagłówków (\textit{title}).
Latencja w~ms na pojedynczą próbkę.}
\label{tab:rek_title}
\begin{tabular}{cllrrrr}
\toprule
\# & Model + emb. & Zbiór tren. & $F_1^{\text{web}}$ & $F_1^{\text{test}}$ & $Q$ & $t_{\text{lat}}$ [ms] \\
\midrule
1 & xgb + tfidf  & WELFake & 0,8499 & 0,9711 & 0,8984 & 0,012 \\
2 & rf  + bow    & ISOT    & 0,8524 & 0,9564 & 0,8940 & 0,009 \\
3 & dt  + bow    & WELFake & 0,8496 & 0,9364 & 0,8843 & 0,007 \\
\bottomrule
\end{tabular}
\end{table}

Wszystkie trzy najlepsze modele dla nagłówków pochodzą z~rodziny
klasycznych algorytmów uczenia maszynowego. \textbf{Rekomendowanym
modelem domyślnym jest XGBoost~+~TF-IDF wytrenowany na WELFake},
łączący wysoką stabilność in-domain (\(F_1^{\text{test}} = 0{,}9711\))
z~dobrą generalizacją na dane WWW (\(F_1^{\text{web}} = 0{,}8499\))
oraz znikomą latencją (\(0{,}012\)~ms na predykcję). W~tym scenariuszu
żaden transformer nie znalazł się w~pierwszej dziesiątce -- modele
BERT, RoBERTa i~DistilBERT trenowane na pełnych artykułach źle
ekstrapolują na krótkie nagłówki.

\subsubsection*{Rekomendacja dla krótkich fragmentów
(\textit{short text}, 30--100~słów)}

\begin{table}[ht]
\centering
\caption{Top~3 modele dla klasyfikacji krótkich fragmentów
(\textit{short text}). Latencja w~ms na pojedynczą próbkę.}
\label{tab:rek_short}
\begin{tabular}{cllrrrr}
\toprule
\# & Model + emb. & Zbiór tren. & $F_1^{\text{web}}$ & $F_1^{\text{test}}$ & $Q$ & $t_{\text{lat}}$ [ms] \\
\midrule
1 & roberta + roberta-base & ISOT    & 0,8556 & 0,9995 & 0,9132 & 0,783 \\
2 & rf + bow               & ISOT    & 0,8747 & 0,9564 & 0,9074 & 0,016 \\
3 & xgb + bow              & WELFake & 0,8349 & 0,9713 & 0,8895 & 0,015 \\
\bottomrule
\end{tabular}
\end{table}

W~scenariuszu krótkich fragmentów najlepszą jakość łączną osiąga
RoBERTa wytrenowana na ISOT, jednak różnica względem klasycznego
RandomForest~+~BoW (również wytrenowanego na ISOT) wynosi zaledwie
\(0{,}58\) punktu miary $Q$, podczas gdy \textbf{Random Forest jest
od transformera około 50 razy szybszy}. Co więcej, RF~+~BoW
odnotowuje \emph{najwyższe} $F_1^{\text{web}}$ spośród wszystkich
badanych modeli w~tym scenariuszu (\(0{,}8747\)). Z~tego powodu jako
\textbf{model domyślny rekomenduje się Random Forest~+~BoW trenowany
na ISOT}, traktując RoBERTa jako opcjonalną ,,wysokojakościową''
ścieżkę uruchamianą np.\ przy niskiej pewności predykcji modelu
szybkiego (kaskadowy mechanizm \textit{confidence fallback}).

\subsubsection*{Rekomendacja dla pełnych artykułów (\textit{text},
$>100$ słów)}

\begin{table}[ht]
\centering
\caption{Top~3 modele dla klasyfikacji pełnych artykułów
(\textit{text}). Latencja w~ms na pojedynczą próbkę.}
\label{tab:rek_text}
\begin{tabular}{cllrrrr}
\toprule
\# & Model + emb. & Zbiór tren. & $F_1^{\text{web}}$ & $F_1^{\text{test}}$ & $Q$ & $t_{\text{lat}}$ [ms] \\
\midrule
1 & bilstm + glove          & LIAR & 0,8550 & 0,6157 & 0,7593 & 0,118 \\
2 & roberta + roberta-base  & ISOT & 0,7724 & 0,9995 & 0,8632 & 0,995 \\
3 & cnn + glove             & LIAR & 0,8244 & 0,6006 & 0,7349 & 0,059 \\
\bottomrule
\end{tabular}
\end{table}

Wybór modelu domyślnego dla pełnych artykułów wymaga osobnego
omówienia, ponieważ kryteria $F_1^{\text{web}}$ i~$F_1^{\text{test}}$
prowadzą do różnych rekomendacji. BiLSTM~+~GloVe trenowany na LIAR
uzyskuje najwyższe $F_1^{\text{web}}$ (\(0{,}8550\)) i~ośmiokrotnie
krótszy czas inferencji niż RoBERTa, lecz jego skuteczność in-domain
jest skromna (\(0{,}6157\)). RoBERTa trenowana na ISOT jest niemal
idealna in-domain (\(0{,}9995\)), ale jej $F_1^{\text{web}}$ jest
o~\(8{,}3\)~pp.\ niższe od najlepszego DL. Z~punktu widzenia wtyczki
Chrome decydujące znaczenie ma $F_1^{\text{web}}$ -- to ona
odzwierciedla rzeczywistą jakość detekcji na artykułach pobieranych
z~Internetu. Dlatego jako \textbf{model domyślny rekomenduje się
BiLSTM~+~GloVe wytrenowany na LIAR}.

\subsubsection*{Podsumowanie rekomendacji}

Tabela~\ref{tab:rek_final} zbiera ostateczne rekomendacje modeli
domyślnych dla mechanizmu \textit{smart routingu} we wtyczce Chrome.

\begin{table}[ht]
\centering
\caption{Ostateczna rekomendacja modeli domyślnych dla wtyczki
Chrome w~zależności od długości tekstu wejściowego.}
\label{tab:rek_final}
\begin{tabular}{lllrr}
\toprule
Zakres długości & Model domyślny & Zbiór tren. & $F_1^{\text{web}}$ & $t_{\text{lat}}$ [ms] \\
\midrule
$<30$ słów        & xgb + tfidf            & WELFake & 0,8499 & 0,012 \\
$30$--$100$ słów  & rf + bow               & ISOT    & 0,8747 & 0,016 \\
$>100$ słów       & bilstm + glove         & LIAR    & 0,8550 & 0,118 \\
\bottomrule
\end{tabular}
\end{table}

Z~tabeli~\ref{tab:rek_final} wynika kilka istotnych obserwacji:

\begin{enumerate}
  \item Dla wszystkich trzech zakresów długości tekstu \emph{latencja
        pozostaje poniżej $0{,}12$~ms na predykcję} -- gwarantuje to
        praktycznie natychmiastową odpowiedź wtyczki nawet przy
        obciążeniu rzędu setek zapytań na sekundę.
  \item Żaden z~rekomendowanych modeli domyślnych nie jest
        transformerem. Mimo że transformery dominują w~ewaluacji
        in-domain, ich przewaga zanika lub odwraca się na danych
        pozyskanych ze stron WWW -- przy znacznie wyższym koszcie
        obliczeniowym.
  \item Dla każdego zakresu długości najlepiej sprawdza się
        \emph{inny zbiór treningowy} (WELFake, ISOT, LIAR), co
        potwierdza tezę o~komplementarności tych korpusów.
  \item Wyniki $F_1^{\text{web}}$ wszystkich trzech rekomendowanych
        modeli mieszczą się w~bardzo wąskim paśmie
        \(0{,}85\)--\(0{,}87\), co świadczy o~spójności jakości
        całego systemu \textit{smart routingu} niezależnie od
        długości tekstu.
\end{enumerate}

\subsection{Implikacje praktyczne}
\label{sec:podsum_praktyczne}

Z~przeprowadzonych eksperymentów wynika kilka rekomendacji
projektowych dla docelowego systemu detekcji fałszywych informacji:

\begin{itemize}
  \item w~scenariuszu rozszerzenia przeglądarki, w~którym istotna
        jest niska latencja, należy preferować klasyczne modele
        z~reprezentacją rzadką;
  \item ze względu na bardzo dobre zachowanie się \textit{RandomForest
        + BoW} na nagłówkach i~lead-ach, ten model został wybrany
        jako klasyfikator domyślny dla krótkich tekstów
        (\textit{smart routing}, sekcja~\ref{sec:podsum_rekomendacja});
  \item dla dłuższych tekstów (powyżej 100~słów) warto rozważyć
        BiLSTM~+~GloVe wytrenowany na LIAR -- jego skuteczność na
        danych spoza zbioru treningowego jest najwyższa w~całym
        badaniu;
  \item transformery (RoBERTa) zaleca się stosować jako klasyfikator
        ,,wysokiej jakości'' uruchamiany na żądanie, np.\ przy
        wykryciu nieoczywistego artykułu, gdzie szybsze modele nie
        są pewne wyniku (niska wartość \textit{confidence});
  \item zbiór LIAR powinien być rozszerzony o~dodatkowe wypowiedzi
        lub zastąpiony nowocześniejszym korpusem -- bez tego
        skuteczność klasyfikacji wypowiedzi politycznych pozostanie
        na poziomie nieakceptowalnym dla zastosowań praktycznych.
\end{itemize}

\subsection{Ograniczenia badania}
\label{sec:podsum_ograniczenia}

Należy zaznaczyć następujące ograniczenia przeprowadzonych
eksperymentów:

\begin{itemize}
  \item liczba prób Optuna była ograniczona (najczęściej do
        $n_\text{trials} = 1$--$10$) ze względu na koszt
        obliczeniowy; nie można więc wykluczyć, że dalsze strojenie
        hiperparametrów poprawiłoby wyniki, zwłaszcza w~rodzinie
        sieci głębokich;
  \item ewaluacja na danych WWW została wykonana na zbiorze
        \(1{,}5\)~tys.\ artykułów i~jest skoncentrowana na
        anglojęzycznych portalach informacyjnych -- uogólnienie do
        innych języków pozostaje otwartym pytaniem badawczym;
  \item modele uczą się cech charakterystycznych dla satyry, a~nie
        intencji manipulacyjnej -- szczegóły uzasadnienia oraz
        konsekwencje zostały omówione w~sekcji~%
        \ref{sec:wyniki_www_dane};
  \item analiza nie uwzględnia kosztu energii ani emisji CO$_2$
        związanego z~uczeniem dużych modeli;
  \item użyte transformery są stosunkowo niewielkie (modele
        \textit{base}); zastosowanie wariantów \textit{large} lub
        modeli typu \textit{instruct-tuned LLM} mogłoby zmienić
        granicę Pareto w~istotny sposób.
\end{itemize}

Pomimo powyższych ograniczeń przeprowadzone badanie pozwoliło na
wszechstronną ocenę różnych podejść do detekcji fałszywych
informacji oraz na sformułowanie konkretnych rekomendacji
projektowych. Zaprezentowane wyniki potwierdzają, że nie istnieje
jeden ,,najlepszy model'' -- wybór konkretnej architektury powinien
być zawsze podejmowany w~kontekście wymagań biznesowych
(czas odpowiedzi, koszt obliczeniowy, oczekiwana stabilność
względem zmian dystrybucji danych).

\chapter{Implementacja praktyczna -- wtyczka oceniająca wiarygodność tekstu}

\section{Założenia projektowe aplikacji}
\label{sec:impl_zalozenia}

Modele wytrenowane w~ramach badań omówionych w~rozdziałach
\ref{sec:wyniki_klasyczne}--\ref{sec:wyniki_kompromis} stanowią
techniczne zaplecze docelowej aplikacji użytkowej: rozszerzenia
do przeglądarki Google~Chrome wspomagającego ocenę wiarygodności
treści napotykanych w~Internecie. Niniejsza sekcja przedstawia
założenia projektowe wynikające z~wymagań funkcjonalnych i~wniosków
z~analizy porównawczej modeli.

\subsection*{Wymagania funkcjonalne}

\begin{itemize}
  \item \textbf{Klasyfikacja binarna.} Aplikacja przyjmuje fragment
        tekstu i~zwraca decyzję (\textit{REAL}/\textit{FAKE}) z~miarą
        pewności z~przedziału $[0,1]$.
  \item \textbf{Praca z~dowolną stroną WWW.} Klasyfikacja odbywa się
        na tekście zaznaczonym przez użytkownika w~aktywnej karcie
        przeglądarki -- niezależnie od domeny, struktury HTML czy języka.
  \item \textbf{Automatyczny dobór modelu.} Użytkownik nie konfiguruje
        żadnych parametrów -- aplikacja samodzielnie wybiera model
        i~interpretuje wynik.
  \item \textbf{Obsługa zbyt krótkiego tekstu.} Dla fragmentów poniżej
        10~znaków aplikacja zwraca komunikat informacyjny zamiast
        klasyfikacji.
\end{itemize}

\subsection*{Wymagania niefunkcjonalne}

\begin{itemize}
  \item \textbf{Niska latencja.} Czas od kliknięcia przycisku do
        wyświetlenia werdyktu nie powinien przekraczać kilkuset
        milisekund. Analiza z~sekcji~\ref{sec:kompromis_pareto}
        potwierdza spełnialność tego wymogu dla rekomendowanych modeli.
  \item \textbf{Prywatność danych.} Analizowany tekst nie opuszcza
        urządzenia użytkownika: komponent obliczeniowy działa lokalnie
        pod adresem \texttt{127.0.0.1}.
  \item \textbf{Modularność.} Logika doboru modelu, obsługa HTTP
        i~warstwa prezentacji są rozdzielone, co umożliwia wymianę
        poszczególnych komponentów bez ingerencji w~pozostałe warstwy.
  \item \textbf{Zgodność ze standardem Manifest~V3.} Wtyczka
        implementuje specyfikację~\cite{chrome_mv3} obowiązującą
        dla wszystkich nowych rozszerzeń Chrome od 2024 roku.
  \item \textbf{Reprodukowalność.} Konfiguracja modeli jest
        zdefiniowana deklaratywnie w~jednym miejscu kodu.
\end{itemize}

\subsection*{Wymagania wynikające z~analizy eksperymentalnej}

Wyniki opisane w~sekcjach \ref{sec:wyniki_www} i~\ref{sec:podsum_rekomendacja}
prowadzą do dwóch dodatkowych założeń:

\begin{enumerate}
  \item \textbf{Wybór modelu zależny od długości tekstu.} Dla każdej
        kategorii długości (\textit{title}, \textit{short text},
        \textit{text}) najlepiej sprawdza się inny model. Aplikacja
        implementuje mechanizm \textit{smart routingu} (sekcja~\ref{sec:impl_router}).
  \item \textbf{Rezygnacja z~transformerów jako domyślnych modeli.}
        Transformery dominują in-domain, lecz ich przewaga zanika
        poza dystrybucją treningową. Dlatego krótkie teksty obsługuje
        szybki klasyczny model (XGBoost~+~TF-IDF), a~długie -- sieć
        BiLSTM~+~GloVe.
\end{enumerate}

\section{Architektura rozwiązania i~przepływ danych}
\label{sec:impl_architektura}

Architektura aplikacji jest dwuwarstwowa i~realizuje wzorzec
\emph{klient--serwer} dostosowany do specyfiki rozszerzenia przeglądarki.
Warstwa kliencka (wtyczka Chrome) obsługuje interakcję z~użytkownikiem
i~pozyskuje tekst z~aktywnej karty, natomiast serwer (lokalna usługa
REST) wykonuje inferencję modeli. Taki podział pozwala utrzymać cienką
warstwę prezentacji w~JavaScripcie, jednocześnie korzystając z~pełnego
ekosystemu Pythona (PyTorch, scikit-learn, transformers) po stronie
obliczeniowej.

\subsection*{Komponent serwerowy}

Serwer zaimplementowano w~\textit{FastAPI}~\cite{fastapi_docs}
uruchamianym pod kontrolą serwera ASGI \textit{Uvicorn}~\cite{uvicorn_docs}.
Framework ten oferuje: natywne wsparcie dla asynchronicznego przetwarzania
(wysoka przepustowość przy minimalnym narzucie), automatyczną walidację
danych wejściowych za pomocą \textit{Pydantic}~\cite{pydantic_docs},
interaktywną dokumentację Swagger UI pod \texttt{/docs} oraz zgodność
z~architekturą REST~\cite{fielding2000rest} ułatwiającą przyszłą rozbudowę.

API udostępnia pojedynczy endpoint \texttt{POST /predict} przyjmujący
ciało żądania \texttt{\{"text": "..."\}}. Modele są ładowane jednorazowo
przy starcie serwera w~kontekście \textit{lifespan} FastAPI i~rezydują
w~pamięci RAM przez cały czas działania procesu, co eliminuje koszt
deserializacji przy każdym zapytaniu.

Komunikację między wtyczką a~serwerem w~różnych \textit{origin}
umożliwiają nagłówki CORS~\cite{w3c_cors}. W~bieżącej konfiguracji
ustawiono \texttt{allow\_origins=["*"]} -- akceptowalne dla serwera
nasłuchującego wyłącznie na pętli zwrotnej \texttt{127.0.0.1}.
W~scenariuszu wdrożenia publicznego listę dozwolonych origin należałoby zawęzić.

\subsection*{Komponent kliencki -- wtyczka Chrome}

Wtyczka implementuje standard \textit{Manifest~V3}~\cite{chrome_mv3}.
Uprawnienia deklarowane są jawnie w~\texttt{manifest.json}: aplikacja
korzysta wyłącznie z~\texttt{activeTab} (dostęp do aktywnej karty)
i~\texttt{scripting} (wstrzykiwanie kodu w~kontekście strony). Nie
używa przestarzałych stron w~tle (\textit{background pages}) -- wszystkie
operacje wykonywane są w~oknie popup i~dynamicznie wstrzykiwanych skryptach.

Warstwa kliencka składa się z~trzech zasobów statycznych:
\texttt{popup.html} (struktura interfejsu), \texttt{styles.css}
(prezentacja i~responsywne style) oraz \texttt{popup.js} (pobieranie
zaznaczonego tekstu, komunikacja z~API, renderowanie wyników).
Ikona \texttt{icon.png} osadzana jest w~pasku narzędzi na podstawie
sekcji \texttt{action} w~manifeście.

\subsection*{Przepływ danych}

Cykl przetwarzania pojedynczego zapytania przebiega następująco:

\begin{enumerate}
  \item Użytkownik zaznacza fragment tekstu i~klika ikonę wtyczki --
        otwiera się okno popup.
  \item Po naciśnięciu \emph{Analyze Selection} skrypt \texttt{popup.js}
        wywołuje \texttt{chrome.scripting.executeScript()}, wstrzykując
        funkcję odczytującą zaznaczenie przez \texttt{window.getSelection()}.
  \item Pobrany tekst jest przesyłany do serwera żądaniem
        \texttt{POST} z~\texttt{Content-Type: application/json}.
  \item Serwer waliduje dane (Pydantic), sprawdza długość tekstu,
        wywołuje heurystykę routingu (sekcja~\ref{sec:impl_router})
        i~metodę \texttt{predict} wybranego modelu.
  \item Wynik z~metadanymi (nazwa modelu, zbiór treningowy, liczba
        słów) jest serializowany do JSON i~zwracany klientowi.
  \item Wtyczka renderuje werdykt: kolor etykiety (\textit{REAL}/\textit{FAKE}),
        wartość procentowa pewności i~pasek postępu.
\end{enumerate}

\begin{figure}[ht]
\centering
\begin{tikzpicture}[
  node distance     = 0.85cm and 1.5cm,
  every node/.style = {font=\small},
  block/.style      = {rectangle, rounded corners=3pt, draw=black!70,
                       line width=0.4pt, minimum height=0.9cm,
                       minimum width=8cm, align=center, inner sep=6pt},
  user/.style       = {block, fill=black!8,    minimum width=4cm},
  client/.style     = {block, fill=blue!8},
  server/.style     = {block, fill=green!8},
  model/.style      = {block, fill=orange!12},
  arrow/.style      = {-{Latex[length=3mm]}, thick, draw=black!60},
  edgelbl/.style    = {font=\footnotesize\itshape, fill=white,
                       inner xsep=3pt, inner ysep=1pt}
]

\node[user]                                 (user)    {Użytkownik};
\node[client,  below=of user]               (popup)
      {\texttt{popup.html} + \texttt{popup.js} \,(Manifest V3) \\
       \texttt{chrome.scripting.executeScript()} +
       \texttt{window.getSelection()}};
\node[server,  below=1.1cm of popup]        (fastapi)
      {\textbf{FastAPI + Uvicorn} \\
       walidacja \texttt{TextRequest} (Pydantic)};
\node[server,  below=of fastapi]            (router)
      {\texttt{select\_best\_model(text)} \\
       heurystyka długości tekstu (\textit{smart routing})};
\node[model,   below=of router]             (predict)
      {\texttt{ModelFactory.get\_model().predict()} \\[2pt]
       XGBoost~+~TF--IDF \;\textbar\; RF~+~BoW \;\textbar\;
       BiLSTM~+~GloVe};
\node[client,  below=1.1cm of predict]      (display)
      {\texttt{displayResult()} \\
       renderowanie werdyktu (\textit{REAL\,/\,FAKE}) w~popup};

\draw[arrow] (user)    -- node[edgelbl, right]
             {zaznaczenie tekstu w~karcie przeglądarki} (popup);
\draw[arrow] (popup)   -- node[edgelbl, right]
             {HTTP POST \texttt{localhost:8000/predict}} (fastapi);
\draw[arrow] (fastapi) -- (router);
\draw[arrow] (router)  -- (predict);
\draw[arrow] (predict) -- node[edgelbl, right]
             {JSON \texttt{\{label, score, meta\}}} (display);

\node[anchor=west, font=\footnotesize\sffamily, text=blue!60!black]
      at ([xshift=4pt]popup.east)   {klient};
\node[anchor=west, font=\footnotesize\sffamily, text=green!50!black]
      at ([xshift=4pt]fastapi.east) {serwer};
\node[anchor=west, font=\footnotesize\sffamily, text=orange!70!black]
      at ([xshift=4pt]predict.east) {modele};
\node[anchor=west, font=\footnotesize\sffamily, text=blue!60!black]
      at ([xshift=4pt]display.east) {klient};

\end{tikzpicture}
\caption{Schematyczny przepływ danych w~aplikacji od zaznaczenia
tekstu przez użytkownika do wyświetlenia werdyktu. Kolorystyka
oznacza warstwę logiczną: niebieski -- wtyczka Chrome (klient),
zielony -- API REST (serwer), pomarańczowy -- inferencja modelu.}
\label{fig:przeplyw_danych}
\end{figure}

Całkowity czas obsługi zapytania (\textit{end-to-end latency}) wynosi
w~praktyce 10--200~ms dla modeli klasycznych i~sieci głębokich
oraz 200--1500~ms dla transformerów, w~zależności od długości tekstu.

\section{Heurystyka doboru modelu w~czasie rzeczywistym}
\label{sec:impl_router}

Kluczowym elementem serwera jest mechanizm \textit{smart routingu}
-- doboru modelu do długości tekstu wejściowego. Heurystyka ta jest
bezpośrednią realizacją rekomendacji z~sekcji~\ref{sec:podsum_rekomendacja}
(Tabela~\ref{tab:rek_final}): dla każdej z~trzech kategorii długości
wybrano model optymalny pod względem $F_1^{\text{web}}$ i~latencji.

\subsection*{Zasada działania}

Serwer wyznacza długość tekstu jako liczbę tokenów oddzielonych
białymi znakami (\texttt{len(text.split())}), a~następnie kieruje
zapytanie do jednego z~trzech predefiniowanych modeli:

\begin{table}[ht]
\centering
\caption{Reguły doboru modelu w~zależności od długości tekstu
wejściowego. Konfiguracja jest spójna z~rekomendacją wyprowadzoną
w~sekcji~\ref{sec:podsum_rekomendacja} (por.\
Tabela~\ref{tab:rek_final}).}
\label{tab:routing}
\small
\setlength{\tabcolsep}{6pt}
\begin{tabular}{@{}llll@{}}
\toprule
Liczba słów & Kategoria (\texttt{api.py}) & Model & Zbiór treningowy \\
\midrule
$<30$        & \texttt{short\_text}   & XGBoost~+~TF--IDF & WELFake \\
$30$--$100$  & \texttt{general}       & RF~+~BoW          & ISOT    \\
$>100$       & \texttt{long\_article} & BiLSTM~+~GloVe    & LIAR    \\
\bottomrule
\end{tabular}
\end{table}

\subsection*{Uzasadnienie wyboru poszczególnych modeli}

Każdy z~trzech modeli został dobrany niezależnie, na podstawie
wyników ewaluacji in-domain oraz na danych pozyskanych ze stron WWW
(sekcje~\ref{sec:wyniki_klasyczne}--\ref{sec:wyniki_www}).
Uzasadnienia są następujące:

\begin{itemize}
  \item \textbf{XGBoost~+~TF-IDF (WELFake) dla nagłówków
        ($<30$~słów).} Model ten łączy wysoką stabilność in-domain
        ($F_1^{\text{test}} = 0{,}9711$) z~bardzo dobrą
        generalizacją na dane WWW
        ($F_1^{\text{web}} = 0{,}8499$) przy znikomej latencji
        rzędu \(0{,}012\)~ms na predykcję. Wybór zbioru WELFake jest
        nieprzypadkowy -- jako korpus zagregowany z~czterech
        podzbiorów cechuje się większą różnorodnością stylistyczną
        niż ISOT, dzięki czemu wytrenowane na nim klasyfikatory
        lepiej znoszą krótkie, jednozdaniowe nagłówki redakcyjne.
        W~rankingu z~Tabeli~\ref{tab:rek_title} model ten zajmuje
        pierwszą pozycję pod względem miary syntetycznej $Q$.
  \item \textbf{RandomForest~+~BoW (ISOT) dla krótkich fragmentów
        ($30$--$100$~słów).} Model uzyskał \emph{najwyższy} wynik
        $F_1^{\text{web}} = 0{,}8747$ spośród wszystkich badanych
        konfiguracji w~scenariuszu \emph{short text}
        (Tabela~\ref{tab:rek_short}), wyprzedzając nawet RoBERTa
        wytrenowaną na tym samym zbiorze. Drzewiasta struktura
        modelu w~połączeniu z~reprezentacją BoW okazuje się
        wyjątkowo odporna na zmiany dystrybucyjne -- model
        wykorzystuje proste cechy częstotliwościowe zamiast
        złożonych semantycznych, co lepiej generalizuje na artykuły
        z~portali spoza zbioru treningowego. Dodatkowym argumentem
        jest niski koszt obliczeniowy (\(0{,}016\)~ms na predykcję)
        oraz brak konieczności udostępniania karty GPU.
  \item \textbf{BiLSTM~+~GloVe (LIAR) dla długich tekstów
        ($>100$~słów).} Model ten osiągnął najwyższy wynik
        $F_1^{\text{web}} = 0{,}8550$ na pełnych tekstach artykułów
        spośród wszystkich badanych konfiguracji
        (sekcja~\ref{sec:wyniki_www_podsum}). Zaskakujący wybór
        zbioru treningowego LIAR wynika z~tego, że krótkie
        wypowiedzi polityczne -- choć trudne in-domain -- zmuszają
        model do nauki ogólnych cech językowych dezinformacji,
        które następnie dobrze przenoszą się na artykuły pełnej
        długości. Latencja rzędu \(0{,}118\)~ms na predykcję jest
        akceptowalna dla scenariusza, w~którym użytkownik
        zaznacza obszerny fragment tekstu i~jest gotów odczekać
        chwilę na wynik.
\end{itemize}

\subsection*{Mechanizm rezerwowy (\textit{fallback})}

Jeśli docelowy model jest niedostępny (brak zapisanych wag lub błąd
inicjalizacji), router przekierowuje zapytanie do modelu \texttt{general}
(RF + BoW na ISOT). Dzięki dobrej generalizacji i~odporności
na zmiany dystrybucyjne jest on w~stanie obsłużyć teksty spoza
swojego naturalnego zakresu długości. Jeżeli i~ten model nie
został załadowany, system korzysta z~pierwszego dostępnego
w~kolekcji; kod HTTP~503 jest zwracany dopiero przy całkowitym
braku modeli.

\subsection*{Uzasadnienie wyboru progów}

Progi 30 i~100 słów wynikają z~typowych rozkładów długości tekstu
w~zbiorze ewaluacyjnym ze stron WWW (sekcja~\ref{sec:wyniki_www_dane}):
30 słów odpowiada granicy między tytułem a~leadem, 100 słów --
między leadem a~pełnym artykułem. Wartości te są decyzją heurystyczną
spójną z~segmentacją stosowaną podczas ewaluacji web, nie zaś wynikiem
algorytmicznej optymalizacji. Możliwym kierunkiem rozwoju jest
zastosowanie \textit{confidence fallback}: model szybki uruchamia
model wolniejszy tylko wówczas, gdy własna pewność predykcji
jest poniżej przyjętego progu.

\section{Interfejs użytkownika i~testy działania}
\label{sec:impl_ui}

Sekcja omawia warstwę prezentacji wtyczki oraz wyniki testów
funkcjonalnych i~wydajnościowych w~rzeczywistych warunkach pracy.

\subsection*{Projekt interfejsu użytkownika}

Interfejs zaprojektowano zgodnie z~zasadą \emph{minimalnej liczby
kroków}: ocena dowolnego fragmentu tekstu wymaga co najwyżej dwóch
kliknięć (ikona wtyczki + \emph{Analyze Selection}).
Okno popup składa się z~trzech elementów wizualnych:

\begin{figure}[ht]
\centering
\includegraphics[width=0.42\textwidth]{img/popup_initial.png}
\caption{Stan początkowy wtyczki -- popup widoczny po kliknięciu
ikony w~pasku narzędzi Chrome. Użytkownik jest informowany
o~konieczności zaznaczenia tekstu na stronie.}
\label{fig:popup_initial}
\end{figure}

\begin{itemize}
  \item \textbf{Nagłówek aplikacji} -- ikona oraz tytuł
        ,,Fake News Detector'' wyraźnie identyfikujące narzędzie.
  \item \textbf{Pole statusu} -- dynamicznie zmieniający się
        komunikat informujący o~bieżącym stanie aplikacji
        (oczekiwanie na zaznaczenie, pobieranie tekstu, analiza,
        komunikat o~zakończeniu).
  \item \textbf{Karta wyniku} -- widoczna po zakończeniu analizy:
        zawiera etykietę (\textit{REAL}~/~\textit{FAKE})
        z~kolorowym wyróżnieniem (zielony / czerwony), miarę
        pewności wyrażoną w~procentach, pasek postępu wizualnie
        odwzorowujący wartość pewności oraz informację o~nazwie
        modelu, który wykonał klasyfikację.
\end{itemize}

Wybór binarnej etykiety i~kolorowego wyróżnienia wynika z~badań
psychologii odbioru dezinformacji. Pennycook i~Rand~\cite{pennycook2021psychology}
wskazują, że proste, jednoznaczne ostrzeżenia są skuteczniejsze
w~modyfikacji oceny wiarygodności niż złożone wskaźniki numeryczne.
Jednocześnie Pennycook i~in.~\cite{pennycook2020implied} ostrzegają
przed \emph{efektem implikowanej prawdy}: brak ostrzeżenia bywa
interpretowany jako potwierdzenie prawdziwości. Dlatego aplikacja
\emph{zawsze} wyświetla werdykt -- niezależnie od klasy
(\textit{REAL} lub \textit{FAKE}) -- zamiast oznaczać wyłącznie
fałszywe treści.

Każda predykcja zawiera też adnotację ,,Analyzed by: \textit{[nazwa modelu]}'',
uświadamiając użytkownikowi, że ocena pochodzi od konkretnego
klasyfikatora i~stanowi wskazówkę, nie rozstrzygające orzeczenie.

\begin{figure}[ht]
\centering
\begin{subfigure}[t]{0.45\textwidth}
  \centering
  \includegraphics[width=\textwidth]{img/popup_real.png}
  \caption{Werdykt \textit{REAL} (artykuł z~BBC News).}
  \label{fig:popup_real}
\end{subfigure}\hfill
\begin{subfigure}[t]{0.45\textwidth}
  \centering
  \includegraphics[width=\textwidth]{img/popup_fake.png}
  \caption{Werdykt \textit{FAKE} (tekst z~The Onion).}
  \label{fig:popup_fake}
\end{subfigure}
\caption{Wyniki klasyfikacji wyświetlane w~popupie. Kolor i~ikona
jednoznacznie identyfikują klasę (zielony dla \textit{REAL},
czerwony dla \textit{FAKE}), miara pewności jest prezentowana
zarówno liczbowo, jak i~wizualnie przez pasek postępu.}
\label{fig:popup_results}
\end{figure}

\begin{figure}[ht]
\centering
\includegraphics[width=0.95\textwidth]{img/extension_in_context.png}
\caption{Wtyczka w~pełnym kontekście użytkowania: zaznaczony
fragment artykułu w~karcie przeglądarki Chrome oraz popup
prezentujący wynik klasyfikacji. Screenshot ilustruje typowy
przepływ pracy użytkownika końcowego.}
\label{fig:extension_context}
\end{figure}

\subsection*{Obsługa przypadków szczególnych}

Wtyczka obsługuje następujące przypadki graniczne:

\begin{itemize}
  \item \textbf{Brak zaznaczenia lub zaznaczenie poniżej 10~znaków}
        -- wyświetlany jest komunikat informacyjny, zapytanie
        do API nie jest wysyłane.
  \item \textbf{Brak połączenia z~serwerem} -- w~razie błędu
        \texttt{fetch} pojawia się komunikat zachęcający do
        uruchomienia \texttt{api.py}.
  \item \textbf{Błąd serwera} (HTTP $\ne 200$) -- wyświetlany
        jest kod błędu, karta wyniku pozostaje ukryta.
  \item \textbf{Tekst zbyt krótki na poziomie API} (poniżej
        10~znaków po normalizacji) -- serwer zwraca etykietę
        \textit{NEUTRAL}; klient interpretuje to jako stan
        informacyjny, nie jako błąd.
\end{itemize}

\subsection*{Testy funkcjonalne}

Weryfikację poprawności działania przeprowadzono w~trzech
scenariuszach testowych, obejmujących artykuły z~portali
informacyjnych i~odpowiadające im teksty satyryczne:

\begin{table}[ht]
\centering
\caption{Scenariusze testowe weryfikujące działanie mechanizmu
\textit{smart routingu}.}
\label{tab:scenariusze}
\small
\setlength{\tabcolsep}{6pt}
\begin{tabular}{@{}p{3.2cm}p{2cm}p{3.4cm}p{3.4cm}@{}}
\toprule
Scenariusz & Liczba słów & Oczekiwany model & Faktyczna kategoria \\
\midrule
Zaznaczenie nagłówka  & 8--20  & XGBoost~+~TF-IDF (WELFake)
                      & \texttt{short\_text} \\
Zaznaczenie lead-u    & 40--80 & RF~+~BoW (ISOT)
                      & \texttt{general} \\
Zaznaczenie pełnego artykułu
                      & 250--600 & BiLSTM~+~GloVe (LIAR)
                      & \texttt{long\_article} \\
\bottomrule
\end{tabular}
\end{table}

\begin{figure}[ht]
\centering
\begin{subfigure}[t]{0.31\textwidth}
  \centering
  \includegraphics[width=\textwidth]{img/routing_title.png}
  \caption{Nagłówek (\textit{$<$30~słów}) -- \texttt{short\_text}
           $\rightarrow$ XGBoost+TF-IDF.}
  \label{fig:routing_title}
\end{subfigure}\hfill
\begin{subfigure}[t]{0.31\textwidth}
  \centering
  \includegraphics[width=\textwidth]{img/routing_lead.png}
  \caption{Lead (\textit{30--100~słów}) -- \texttt{general}
           $\rightarrow$ RandomForest+BoW.}
  \label{fig:routing_lead}
\end{subfigure}\hfill
\begin{subfigure}[t]{0.31\textwidth}
  \centering
  \includegraphics[width=\textwidth]{img/routing_full.png}
  \caption{Pełny artykuł (\textit{$>$100~słów}) --
           \texttt{long\_article} $\rightarrow$ BiLSTM+GloVe.}
  \label{fig:routing_full}
\end{subfigure}
\caption{Praktyczna weryfikacja mechanizmu \textit{smart routingu}.
Trzy zaznaczenia tekstów o~różnej długości w~tym samym artykule
zostały rozpoznane przez router jako różne kategorie i~każde
zostało obsłużone przez właściwy model -- co jest widoczne w~polu
\texttt{Analyzed by} na karcie wyniku.}
\label{fig:routing_demo}
\end{figure}

We wszystkich trzech scenariuszach mechanizm routingu działał
zgodnie z~oczekiwaniami: pole ,,Analyzed by'' wyświetlało
identyfikator modelu właściwy dla zakresu długości tekstu.
Obsługa przypadków granicznych (zbyt krótkie zaznaczenie, brak serwera)
również przebiegła zgodnie ze specyfikacją.

\subsection*{Testy wydajnościowe}

Czasy odpowiedzi (\textit{end-to-end latency}) zmierzono na maszynie
z~procesorem klasy laptop i~zintegrowaną kartą graficzną (inferencja
wyłącznie na CPU). Wyniki zestawiono w~Tabeli~\ref{tab:latency_e2e}.

\begin{table}[ht]
\centering
\caption{Średnie czasy odpowiedzi aplikacji w~typowych warunkach
pracy. Obejmują pobranie zaznaczenia, transport HTTP, inferencję
oraz renderowanie wyniku.}
\label{tab:latency_e2e}
\small
\setlength{\tabcolsep}{6pt}
\begin{tabular}{@{}lrr@{}}
\toprule
Model (kategoria)              & Średni czas [ms] & Mediana [ms] \\
\midrule
XGBoost~+~TF-IDF (\emph{short\_text}) &  30 &  25 \\
RF~+~BoW         (\emph{general})     &  35 &  28 \\
BiLSTM~+~GloVe   (\emph{long\_article})
                                      & 180 & 160 \\
\bottomrule
\end{tabular}
\end{table}

Wszystkie trzy modele rezydujące w~pamięci oferują średni czas
odpowiedzi poniżej 200~ms, co czyni aplikację praktycznie
nieodczuwalną dla użytkownika. Wyniki są spójne z~analizą z~sekcji~\ref{sec:kompromis_srednie}
i~potwierdzają zasadność wyboru modeli klasycznych i~lekkiej sieci
BiLSTM zamiast kosztownych transformerów.

\subsection*{Podsumowanie sekcji}

Rozszerzenie do przeglądarki Chrome stanowi praktyczną realizację
wniosków z~rozdziałów \ref{sec:wyniki_klasyczne}--\ref{sec:wyniki_kompromis}:
różnym kategoriom długości tekstu odpowiadają wcześniej wyselekcjonowane
modele uczenia maszynowego, a~komunikacja klient--serwer odbywa się
wyłącznie lokalnie. Połączenie nowoczesnych standardów technicznych
(Manifest~V3, FastAPI, Pydantic, REST), decyzji projektowych
opartych na literaturze psychologicznej~\cite{pennycook2021psychology,
pennycook2020implied} i~wniosków empirycznych z~eksperymentów daje
narzędzie, które może realnie wspierać użytkownika w~ocenie wiarygodności
treści napotykanych w~codziennym korzystaniu z~sieci.

\chapter{Podsumowanie i wnioski końcowe}

\section{Ocena stopnia realizacji celu pracy}

\section{Odpowiedzi na pytania badawcze}

\section{Weryfikacja hipotezy}

\section{Wkład badawczy w obszar detekcji fake newsów}

\section{Ograniczenia przeprowadzonych badan i zaimplementowanego narzedzia}

\section{Kierunki dalszego rozwoju}

\printbibliography

\end{document}
